\documentclass[11pt]{article}
\usepackage[margin=1.1in]{geometry}
\usepackage{amsmath,amssymb,amsthm,mathtools,booktabs,array,tabularx,longtable}
\usepackage[dvipsnames]{xcolor}
\usepackage[colorlinks=true,linkcolor=MidnightBlue,citecolor=MidnightBlue,urlcolor=MidnightBlue]{hyperref}
\usepackage{fancyhdr}
\pagestyle{fancy}\fancyhf{}\renewcommand{\headrulewidth}{0pt}
\fancyfoot[R]{\thepage}
\fancyfoot[L]{\scriptsize\color{gray}Generated by the Transport Lab, not the version of record. Official version: \href{https://doi.org/10.5281/zenodo.22722121}{doi:10.5281/zenodo.22722121}}
\fancypagestyle{plain}{\fancyhf{}\fancyfoot[R]{\thepage}\fancyfoot[L]{\scriptsize\color{gray}Generated by the Transport Lab, not the version of record. Official version: \href{https://doi.org/10.5281/zenodo.22722121}{doi:10.5281/zenodo.22722121}}\renewcommand{\headrulewidth}{0pt}}
\newtheorem{theorem}{Theorem}[section]
\newtheorem{proposition}[theorem]{Proposition}
\newtheorem{lemma}[theorem]{Lemma}
\newtheorem{corollary}[theorem]{Corollary}
\newtheorem{conjecture}[theorem]{Conjecture}
\newtheorem{refutation}[theorem]{Refutation}
\newtheorem{prediction}[theorem]{Prediction}
\newtheorem{principle}[theorem]{Principle}
\newtheorem{observation}[theorem]{Observation}
\newtheorem{axiom}[theorem]{Axiom}
\theoremstyle{definition}
\newtheorem{definition}[theorem]{Definition}
\newtheorem{fact}[theorem]{Fact}
\newtheorem{assembled}[theorem]{Assembly}
\newtheorem{postulate}[theorem]{Postulate}
\newtheorem{premise}[theorem]{Premise}
\newtheorem{hypothesis}[theorem]{Hypothesis}
\newtheorem{remark}[theorem]{Remark}
\newtheorem{example}[theorem]{Example}
\theoremstyle{plain}
\setlength{\marginparwidth}{0.8in}
\newcommand{\changedmark}[2]{\marginpar{\raggedright\scriptsize\color{#1}#2}}
\providecommand{\Aa}{}\renewcommand{\Aa}{\mathcal{A}}
\providecommand{\Ad}{}\renewcommand{\Ad}{\widetilde{\mathrm{Ad}}}
\providecommand{\Aev}{}\renewcommand{\Aev}{\mathcal{A}^{+}}
\providecommand{\Aod}{}\renewcommand{\Aod}{\mathcal{A}^{-}}
\providecommand{\B}{}\renewcommand{\B}{\mathrm{B}}
\providecommand{\C}{}\renewcommand{\C}{\mathbb{C}}
\providecommand{\CAR}{}\renewcommand{\CAR}{\mathrm{CAR}}
\providecommand{\CS}{}\renewcommand{\CS}{{\mathbb{S}^{+}}}
\providecommand{\CSm}{}\renewcommand{\CSm}{{\mathbb{S}^{-}}}
\providecommand{\Cl}{}\renewcommand{\Cl}{\mathrm{Cl}}
\providecommand{\Con}{}\renewcommand{\Con}{\mathcal{C}}
\providecommand{\Dl}{}\renewcommand{\Dl}{\Delta}
\providecommand{\EN}{}\renewcommand{\EN}{E_{N}}
\providecommand{\End}{}\renewcommand{\End}{\mathrm{End}}
\providecommand{\F}{}\renewcommand{\F}{\mathbb{F}}
\providecommand{\G}{}\renewcommand{\G}{\Gamma}
\providecommand{\HH}{}\renewcommand{\HH}{\mathbb{H}}
\providecommand{\IS}{}\renewcommand{\IS}{I\sigma_3}
\providecommand{\J}{}\renewcommand{\J}{\mathbf{J}}
\providecommand{\LH}{}\renewcommand{\LH}{L_{H}}
\providecommand{\Mat}{}\renewcommand{\Mat}{\mathrm{M}}
\providecommand{\Mo}{}\renewcommand{\Mo}{M^{\circ}}
\providecommand{\Mp}{}\renewcommand{\Mp}{M^{+}}
\providecommand{\Msun}{}\renewcommand{\Msun}{M_{\odot}}
\providecommand{\Mt}{}\renewcommand{\Mt}{\widetilde{M}}
\providecommand{\Ncone}{}\renewcommand{\Ncone}{\mathcal{N}}
\providecommand{\Neff}{}\renewcommand{\Neff}{N_{\mathrm{eff}}}
\providecommand{\PCS}{}\renewcommand{\PCS}{\mathbb{P}{\mathbb{S}^{+}}}
\providecommand{\PT}{}\renewcommand{\PT}{\mathbb{PT}}
\providecommand{\Pin}{}\renewcommand{\Pin}{\mathrm{Pin}}
\providecommand{\Q}{}\renewcommand{\Q}{\mathbb{Q}}
\providecommand{\R}{}\renewcommand{\R}{\mathbb{R}}
\providecommand{\Ric}{}\renewcommand{\Ric}{\mathrm{Ric}}
\providecommand{\SBH}{}\renewcommand{\SBH}{S_{\mathrm{BH}}}
\providecommand{\Sc}{}\renewcommand{\Sc}{\mathbb{S}}
\providecommand{\Sch}{}\renewcommand{\Sch}{\mathrm{P}}
\providecommand{\Sig}{}\renewcommand{\Sig}{\Sigma}
\providecommand{\Som}{}\renewcommand{\Som}{S_\omega}
\providecommand{\Sp}{}\renewcommand{\Sp}{\mathbb{S}}
\providecommand{\Sph}{}\renewcommand{\Sph}{\mathrm{S}}
\providecommand{\Spi}{}\renewcommand{\Spi}{S_\pi}
\providecommand{\Spin}{}\renewcommand{\Spin}{\mathrm{Spin}}
\providecommand{\Sth}{}\renewcommand{\Sth}{S_{\mathrm{th}}}
\providecommand{\T}{}\renewcommand{\T}{\mathbb{T}}
\providecommand{\TH}{}\renewcommand{\TH}{T_{\mathrm{H}}}
\providecommand{\Tr}{}\renewcommand{\Tr}{\operatorname{Tr}}
\providecommand{\UHF}{}\renewcommand{\UHF}{\mathrm{UHF}}
\providecommand{\Z}{}\renewcommand{\Z}{\mathbb{Z}}
\providecommand{\Zc}{}\renewcommand{\Zc}{\mathrm{Z}}
\providecommand{\cmg}{}\renewcommand{\cmg}{\mathrm{cm^{2}\,g^{-1}}}
\providecommand{\dS}{}\renewcommand{\dS}{\mathrm{dS}}
\providecommand{\dsc}{}\renewcommand{\dsc}{\mathrm{dsc}}
\providecommand{\dtPG}{}\renewcommand{\dtPG}{dt}
\providecommand{\dtS}{}\renewcommand{\dtS}{dt_{\mathrm{S}}}
\providecommand{\ellv}{}\renewcommand{\ellv}{\ell}
\providecommand{\er}{}\renewcommand{\er}{e_r}
\providecommand{\et}{}\renewcommand{\et}{e_t}
\providecommand{\ev}{}\renewcommand{\ev}{\mathrm{ev}}
\providecommand{\gl}{}\renewcommand{\gl}{\mathfrak{gl}}
\providecommand{\hatV}{}\renewcommand{\hatV}{\widehat{V}}
\providecommand{\hb}{}\renewcommand{\hb}{\bar{h}}
\providecommand{\hf}{}\renewcommand{\hf}{\hat{f}}
\providecommand{\id}{}\renewcommand{\id}{\mathrm{id}}
\providecommand{\kpc}{}\renewcommand{\kpc}{\mathrm{kpc}}
\providecommand{\lP}{}\renewcommand{\lP}{\ell_{P}}
\providecommand{\mD}{}\renewcommand{\mD}{\mathbf{m}_{D}}
\providecommand{\mb}{}\renewcommand{\mb}{\bar m}
\providecommand{\mnu}{}\renewcommand{\mnu}{\Sigma m_{\nu}}
\providecommand{\nInf}{}\renewcommand{\nInf}{n_\infty}
\providecommand{\nO}{}\renewcommand{\nO}{n_0}
\providecommand{\ninf}{}\renewcommand{\ninf}{n_{\infty}}
\providecommand{\nz}{}\renewcommand{\nz}{n_{0}}
\providecommand{\proj}{}\renewcommand{\proj}{\operatorname{proj}}
\providecommand{\pv}{}\renewcommand{\pv}{\pi}
\providecommand{\rad}{}\renewcommand{\rad}{\mathrm{rad}}
\providecommand{\sgn}{}\renewcommand{\sgn}{\operatorname{sgn}}
\providecommand{\shA}{}\renewcommand{\shA}{\mathcal{A}}
\providecommand{\shC}{}\renewcommand{\shC}{\mathcal{C}}
\providecommand{\shE}{}\renewcommand{\shE}{\mathcal{E}}
\providecommand{\shJ}{}\renewcommand{\shJ}{\mathcal{J}}
\providecommand{\shK}{}\renewcommand{\shK}{\mathcal{K}}
\providecommand{\shL}{}\renewcommand{\shL}{L}
\providecommand{\shW}{}\renewcommand{\shW}{\mathcal{W}}
\providecommand{\sll}{}\renewcommand{\sll}{\mathfrak{sl}}
\providecommand{\so}{}\renewcommand{\so}{\mathfrak{so}}
\providecommand{\su}{}\renewcommand{\su}{\mathfrak{su}}
\providecommand{\supp}{}\renewcommand{\supp}{\operatorname{supp}}
\providecommand{\tr}{}\renewcommand{\tr}{\operatorname{tr}}
\providecommand{\uu}{}\renewcommand{\uu}{\mathfrak{u}}
\providecommand{\br}{}\renewcommand{\br}[1]{\langle #1\rangle}
\providecommand{\bra}{}\renewcommand{\bra}[1]{\langle #1|}
\providecommand{\file}{}\renewcommand{\file}[1]{\texttt{#1}}
\providecommand{\hp}{}\renewcommand{\hp}[2]{( #1 \,|\, #2 )}
\providecommand{\ip}{}\renewcommand{\ip}[2]{\langle #1, #2\rangle}
\providecommand{\ket}{}\renewcommand{\ket}[1]{|#1\rangle}
\providecommand{\rel}{}\renewcommand{\rel}[2]{D\bigl(#1 \,\|\, #2\bigr)}
\providecommand{\rev}{}\renewcommand{\rev}[1]{\widetilde{#1}}
\providecommand{\slashed}{}\renewcommand{\slashed}[1]{{#1}\mkern-10mu/}
\title{\bfseries Hawking Radiation from the Ledger\\[4pt]\large\itshape State, Continuum, Lift, and the Finite-Ledger Seam}
\author{Leon Butler\\ \normalsize\href{mailto:lbutler.math@gmail.com}{\texttt{lbutler.math@gmail.com}}\\ \normalsize ORCID: \href{https://orcid.org/0009-0006-4449-5617}{0009-0006-4449-5617}}
\date{Generated by the Transport Lab, 27 October 2026 \textperiodcentered{} Transport Papers, T7\\[3pt]\small Official version: \href{https://doi.org/10.5281/zenodo.22722121}{doi:10.5281/zenodo.22722121}}
\begin{document}
\maketitle
\begin{abstract}
\normalsize
This paper is a quantum companion to T5: it supplies the state, the continuum, and the four-dimensional lift between T5's single-mode core and the physical Hawking spectrum, each claim at its tier. Tier A: among the quasi-free states of the boundary modes' CAR algebra exactly one is KMS for the geometry's boost at the Hawking temperature, the I-state; existence is window-exact, carried to the full algebra by imported inductive-limit results; the KMS identity itself forces the pairing to vanish and the symbol to be Fermi-Dirac, all modes at once, extending T5's single-mode occupancy theorem; uniqueness beyond the class is the imported ideal-Fermi-gas theorem; and the statements carry to the continuum on a dense core. Tier B, the joints itemised: the helicity lift and greybody inheritance, whose only joint proved in this paper, the homogeneity $-2s-2$, is promoted to a stated result from inside the proof of the twistor paper's dictionary theorem; the exterior coincides with general relativity on the invertible region at $\Lambda = 0$, so the flux prediction equals Hawking's, filtering included. Tier A again: the finite mode space's deviations from the free-boson current algebra are edge-supported by an exact commutator identity; and, on the linear s-wave ladder of T5's premise (Q), an exact Planck factor with proven exponential window tails. The $\ell_P^2/A$ scale of generic corrections is a dimensional reading at modelling level; the per-observable scaling is open. Open and named: axiomatic selection of the dynamics, the lift beyond assembly, the Kerr continuum. Exact-arithmetic checks verify the identities.
\end{abstract}
\medskip
\noindent\textbf{MSC 2020.} Primary 81T05, 46L60; Secondary 83C57, 46L30, 81R10.\\
\textbf{Keywords.} Hawking radiation; KMS states; quasi-free states; CAR algebra; modular theory; thermofield double; Penrose transform; greybody factors; bosonisation.
\newpage
\noindent \textbf{Authorship and disclosure. }The founding idea of the series is the author's: that division by zero in a Clifford algebra should be read not as a prohibition but as a transport into a higher-dimensional algebra, and that the transported division need not be accompanied by an inverse operation. The direction of investigation and the critical review of the results are likewise the author's. The mathematical development and the drafting of the text were carried out in extended collaboration with Claude Fable 5 (Anthropic), an AI system.

\noindent\fbox{\parbox{\dimexpr\linewidth-2\fboxsep-2\fboxrule}{\textbf{This is not the version of record.} The official version of this paper is the one deposited on Zenodo, \href{https://doi.org/10.5281/zenodo.22722121}{doi:10.5281/zenodo.22722121}. This PDF was generated by the Transport Lab on 27 October 2026 and assembled from approved parts of its store: the results, proofs and text of the paper as the store holds them, each result with its premises and machine checks. It is regenerated when a result it contains is corrected, so it may differ from the deposited paper. Since the deposited version 1.0: 1 result corrected, 2 changed, 9 reworded (wording, notation or citations only) and 26 cited results added. Each revised or added result is marked in the margin (\textcolor{BrickRed}{corrected}, \textcolor{BurntOrange}{changed}, \textcolor{gray}{reworded}), and the last section, \emph{Changes since the deposited version}, lists them with the reason for each.}}
\newpage

\section{Introduction}

\phantomsection\label{c:note.hawking.setting-introduction}
T5 of this series \cite{ButlerT5} computes a six-step chain from the Schwarzschild gauge field to the Page curve. Two of its links are deliberately minimal, and say so: single-mode thermality is proved from two named premises, (P) univalence and (R) the KMS condition at the geometric period, with the occupancy derived from the gluing identity \cite[Prems.~7.8 and 7.19, Thm.~7.20]{ButlerT5}; and the bosonic Planck sector is derived in the s-wave-reduced theory under a joint quasi-freeness premise \cite[Prem.~8.1]{ButlerT5}. Between that core and the physical Hawking spectrum stand four questions, and this paper answers them at their correct levels.

\emph{Which state?} the section ``The state input: the I-state'' proves the I-state theorem, the I-state being the series' name for the framework's distinguished equilibrium state $\omega_I$: among the quasi-free states of the CAR algebra of the boundary modes there is exactly one that is KMS for the geometry's boost at the Hawking temperature, and the KMS identity forces its two-point data outright, the pairing to zero and the symbol to Fermi-Dirac (Theorem \ref{c:hawking.i-state}). T5 derives the single-mode occupancy the same way; the increment here is operator-level and simultaneous across all modes, with uniqueness inside the quasi-free class proved, and uniqueness beyond it imported with provenance. What is genuinely open is the \emph{selection} of the dynamics from a geometric axiom, and the section ``The state input: the I-state'' poses that open problem precisely, in the terms T5's own selection analysis fixes.

\emph{Which algebra carries the mode continuum?} the section ``The mode continuum'' extends the single-mode statements over the tower's CAR limit \cite{ButlerT3}: the finite-window statements are exact theorems, the inductive-limit results are imported as Facts with provenance, and the result is free-fermion thermal field theory as the algebra carrying the near-horizon boundary modes.

\emph{How does the two-dimensional core become four-dimensional radiation?} the section ``The four-dimensional lift'' assembles the lift with its joints itemised: helicity is the charge of the twistor storey (one level of the completion tower), and the Penrose homogeneity $-2s-2$ is proved from the ordering shift of the twistor paper \cite[Thm.~5.7]{ButlerT9} in Proposition \ref{c:hawking.homogeneity-ordering-shift} (the computation appears inside that theorem's proof; the proposition promotes it to a stated result); greybody factors are inherited verbatim, because the series' exterior description coincides with general relativity on the invertible region at $\Lambda = 0$ \cite{ButlerT6}, so the flux prediction equals Hawking's, filtering included.

\emph{What does the finite boundary mode space change?} the section ``The finite boundary mode space, and the $\ell_P^2/A$ correction'' proves that the deviations of the finite boundary mode space from the infinite current algebra are edge-supported, with the exact edge identity and quantified tails: the corrections are area-suppressed at their proven rates, and the resulting $\ell_P^2/A$ correction is stated as a prediction. The Weinberg-Witten triage against the composite-graviton reading is recorded there as this paper's own claim.

\subsection*{Prior art}
\phantomsection\label{c:note.hawking.setting-prior-art}
The prior work is acknowledged once, here. Hawking's radiance is \cite{Hawking1975}. The thermofield-double reading of the eternal black hole is Israel's \cite{Israel1976}; the identification of modular flow with the boost for wedge algebras is Bisognano-Wichmann \cite{BisognanoWichmann1976}, with Sewell's black-hole reading \cite{Sewell1982}. Quasi-free states of the CAR algebra and their two-point calculus are Araki's \cite{Araki1971}; KMS theory, the ideal Fermi gas, and the inductive-limit results are standard in Bratteli-Robinson \cite{BratteliRobinson1997}. The geometric-modular-action programme, the candidate axiom for state selection, is Buchholz-Dreyer-Florig-Summers \cite{BDFS2000}. The Penrose transform is \cite{Penrose1969, EPW1981}; greybody filtering through the Regge-Wheeler and Teukolsky barriers is \cite{ReggeWheeler1957, Teukolsky1973, Page1976}. The boson-fermion correspondence and the charge-zero Fock isomorphism are classical two-dimensional field theory \cite{FrenkelKac1980}. Weinberg-Witten is \cite{WeinbergWitten1980} and the soft-graviton universality is \cite{Weinberg1964}. What this paper claims as new is the transport assembly: the I-state theorem at the operator level, all modes at once, with pairing and symbol both forced by the KMS identity (the single-mode occupancy is derived the same way in \cite{ButlerT5}); the exact finite-window continuum statements and the limit with its imports declared; the promotion of the homogeneity computation inside the proof of \cite{ButlerT9}'s dictionary theorem to a stated result the transform can consume; the identification of the lift's remaining joints with theorems of this series (the exterior exactness of \cite{ButlerT6}); and the quantified edge-support of the finite-mode corrections.

\subsection*{Contribution}
\phantomsection\label{c:note.hawking.setting-contribution}
At the standard of \cite{ButlerT1}, with imports as Facts and the assembly joints itemised where a result is assembled.

\begin{enumerate}
\item The single-mode modular structure reproved self-contained, and the I-state theorem: existence of the quasi-free state whose modular clock is the geometry's clock, with gauge invariance and the Fermi-Dirac symbol both forced by the KMS identity, extending T5's single-mode derivation \cite[Thm.~7.20]{ButlerT5} to the full algebra, all modes at once; uniqueness beyond the quasi-free class as a corollary of the imported ideal-Fermi-gas theorem; the selection problem posed against T5's own analysis (the section ``The state input: the I-state'').
\item The continuum: exact finite-window KMS theorems, the inductive-limit assembly with its imports stated as Facts, free-fermion thermal field theory (the section ``The mode continuum'', Remark \ref{c:note.hawking.free-fermion-thermal-field-theory-israels}).
\item The lift: the Penrose homogeneity $-2s-2$ stated and proved from the twistor paper's ordering shift (Proposition \ref{c:hawking.homogeneity-ordering-shift}; the computation appears inside the proof of \cite[Thm.~5.7]{ButlerT9}, and is promoted here to a result the transform can consume), discharging the homogeneity half of the successor problem the twistor paper assigns to this paper, the transform itself staying open; the helicity lift and greybody inheritance assembled with every joint flagged, the interface premise (I) named in the statement (the section ``The four-dimensional lift'').
\item The $\ell_P^2/A$ correction: the exact edge identity for the finite window's current commutator, quantified tails, the area-suppressed reading, and the Weinberg-Witten triage (the section ``The finite boundary mode space, and the $\ell_P^2/A$ correction'').
\item Machine verification of the computational identities in exact arithmetic \cite{ButlerVerifyT7}.
\end{enumerate}

\section*{What we assume}
The results below are taken as given in this paper; each result that rests on one says so.
\begin{postulate}[Joint J3: the s-wave reduction (modelling)]\label{c:hawking.joint-j3-s-wave-reduction}
The current boson of the s-wave theory is the s-wave reduction of each four-dimensional massless field, and the lift from the boundary modes to the field is the standard partial-wave reconstruction: the near-horizon reduction of field theory on the Schwarzschild exterior identifies each $(\omega, \ell, m, s)$ channel with a two-dimensional mode.
\end{postulate}

\begin{postulate}[Premise (I): the interface identification]\label{c:hawking.premise-i-interface-identification}
The horizon-sector state feeds the exterior modes in the standard way, i.e.\ the occupancy of each near-horizon channel is carried to the corresponding asymptotic channel by the exterior propagation.
\end{postulate}

\begin{postulate}[(P): univalence]\label{c:singularity.p-univalence}
Physical states of the boundary modes are even: $\omega \circ \alpha = \omega$, equivalently $\omega$ annihilates the odd sector $\Aod$.
\end{postulate}

\begin{postulate}[(Q): joint quasi-freeness, with the s-wave weights at the half-integer offset]\label{c:singularity.q-joint-quasi-freeness-s-wave}
In the stationary idealisation of the boundary modes, and for this step alone, the joint state of the mode array is taken to be the gauge-invariant quasi-free state with the marginal data that (P) and (R) (Premise \ref{c:singularity.r-regularity-fixed-locus}) fix unconditionally (Theorem \ref{c:singularity.rate-discharge-hypothesis-iii}, Corollary \emph{The chain's state class, and its stability} [T5, Cor 7.32]): writing $\psi_k$ for the fermion modes of the array of Theorem \ref{c:singularity.bosonisation-bridge-s-wave-sector}, its two-point function is $\langle \psi^\dagger_k \psi_l \rangle = \delta_{kl}\,n_k$, and every higher correlation function factorises by Wick's theorem, so that all connected correlators vanish and the joint state is KMS at $\beta = 2\pi/\kappa$ for the joint modular flow. The mode weights are evenly spaced with the array's sea level midway between two sites,
\[
\omega_k = \kappa\bigl(k + \tfrac12\bigr), \qquad k \in \Z ,
\]
so that no mode has zero weight, the weights are symmetric under the particle-hole relabelling $k \mapsto -1-k$, and the Killing flow at one modular period $2\pi/\kappa$ is the parity automorphism $\psi \mapsto -\psi$.

\emph{Why this offset.} The weights $\omega_k$ of Definition \emph{The geometric modular condition (G)} [T5, Def 7.12] are fixed nowhere else in the framework: given (P), (R) and the quasi-free clause, an affine law $\omega_k = ck + d$ is what the currents' being normal modes of the boost flow requires, particle-hole symmetry leaves exactly two offset classes, $d = 0$ and $d = c/2$, and the slope is the value $\beta c = 2\pi$ (Theorem \emph{Premise (Q) splits into the maximum-entropy member of the (P),(R) class and an affine weight law with two free data}(ii), (iv)). The class $d = c/2$ is adopted as a consequence of one reading, not as a bare premise: the array's modes are read as the discretised modes of the field theory's spinor fields on the horizon's light rays, $\psi(u) \mapsto e^{\kappa t/2}\psi(e^{\kappa t}u)$, of boost weight $\tfrac12$ (the weight at which the boost preserves the anticommutator), in the state that Premise \emph{(N$_{+}$): on the continuum of a horizon's null generators the state of the modes is a gauge-invariant quasi-free ground state for the affine null translations along the generators} (N$_{+}$) selects, a ground state for the affine null translations of the generators (Theorem \emph{A ground state for the affine null translations of a horizon's generators restricts on either side of the edge to the continuum I-state at rapidity period $2\pi$, so the continuum restatements of (R) for the eternal hole, (R$_{\Lambda}$), (R$_{\mathrm{loc}}$) and the joint KMS part of (Q) follow from one premise with no temperature in it, posable only on the continuum}). A field of boost weight $\tfrac12$ changes sign under a full turn of the regular Euclidean section, and $\kappa(k + \tfrac12)$ is the offset at which the array's Killing flow at one modular period carries that sign (Theorem \emph{With the sea level between sites the s-wave array carries exactly half of ln 2 at the Hawking temperature}). The reading is a decision, not a derivation: the Euclidean-period sign $R_{2\pi} = -1$ (Proposition \emph{Where the $+1$ of Fermi-Dirac comes from} [T5, Cor 5.4]) holds in a KMS state of the array at either offset, so identifying the array's real-time recurrence at one period with it is part of the reading.

\emph{On the continuum.} The joint KMS part of this premise is not a separate assumption on the continuum mode space of a horizon's generators: there it follows from (N$_{+}$), whose null ground state restricts on either side of the edge to the continuum I-state at $\beta = 2\pi/\kappa$, jointly KMS and even (Theorem \emph{A ground state for the affine null translations of a horizon's generators restricts on either side of the edge to the continuum I-state at rapidity period $2\pi$, so the continuum restatements of (R) for the eternal hole, (R$_{\Lambda}$), (R$_{\mathrm{loc}}$) and the joint KMS part of (Q) follow from one premise with no temperature in it, posable only on the continuum}(ii), (iii)). The stationary array statement here is its point-spectrum discretisation, on which (N$_{+}$) itself cannot be posed (part (iv) there). The quasi-free clause is a hypothesis of (N$_{+}$) as well, and (N$_{+}$) does not give the weight clause.

\emph{Consequences of the offset.} At these weights one array carries exactly $\tfrac12\ln 2$ at the Hawking temperature, given the classical theta-function identities, so under the entropy-matching hypothesis two arrays carry the capacity $\ln 2$ of one registered boundary mode with no zero-weight mode (Theorem \emph{With the sea level between sites the s-wave array carries exactly half of ln 2 at the Hawking temperature}(i), (iii)); the current sector of one array is one oscillator tower, one helicity channel, at either offset (Theorem \emph{Under (Q) an array's ln 2 is its zero-mode charge doublet, which no current with m != 0 connects, and its current sector is one oscillator tower}(ii), (iv)), so one registered mode then corresponds to two arrays and two helicity channels. Since every mode has nonzero weight, the gluing equation of Theorem \ref{c:singularity.rate-discharge-hypothesis-iii} fixes $\beta$ at every mode of the array, and the zero-weight scopes of Theorems \ref{c:singularity.descended-form-state-independent}, \emph{Universality of the temperature under (G)} [T5, Cor 7.15] and \ref{c:singularity.rate-discharge-hypothesis-iii} do not arise on it.

\emph{Version 2.} The source states the weight clause as $\omega_k = \kappa k$ with the sea level at $k = 0$, one mode of zero weight at which the gluing equation fixes $n_0 = \tfrac12$ but not $\beta$. Results computed under that offset refer to version 2.
\end{postulate}

\begin{postulate}[(R): regularity at the fixed locus]\label{c:singularity.r-regularity-fixed-locus}
For every mode $k$ and every pair $a, b$ of elements of its ladder $\{v_k, v_k^{*}\}$, the correlation function $F_{a,b}(t) := \omega\bigl(a\,\alpha_t(b)\bigr)$ of the Killing flow $\alpha_t = \operatorname{Ad}(e^{it\omega_k K_k})$, an entire function of $t$, continues in imaginary Killing time to a single-valued function on the regular Euclidean section of the Schwarzschild gravitational field on flat spacetime (Proposition \ref{c:singularity.euclidean-section-period}; Definition \emph{Gauge Theory Gravity: the gravitational field on flat spacetime}): its two orderings are the two sides of the cut, and they meet after one period,
\[
F_{a,b}\bigl(t + i\,p\bigr) \;=\; \omega\bigl(\alpha_t(b)\,a\bigr), \qquad p = \frac{2\pi}{\kappa} .
\]
In words: each single-mode marginal is KMS for the Killing flow at exactly the period the Euclidean section's smoothness at its tip dictates. The regularity is local, posed at the tip of the Euclidean section (a polar chart about the tip of the Wick-rotated flat background); for a state that is KMS for the Killing flow, this is also the period at which it is regular across the horizon locus in the Painlev\'e-Gullstrand chart.

\emph{Note (continuum restatement).} Posed on the continuum of the horizon's null generators (Definition \ref{c:obj.hawking.continuum-mode-space}), the content of (R) for the eternal (Kruskal) hole follows from Premise \emph{(N$_{+}$): on the continuum of a horizon's null generators the state of the modes is a gauge-invariant quasi-free ground state for the affine null translations along the generators}: by Theorem \emph{A ground state for the affine null translations of a horizon's generators restricts on either side of the edge to the continuum I-state at rapidity period $2\pi$, so the continuum restatements of (R) for the eternal hole, (R$_{\Lambda}$), (R$_{\mathrm{loc}}$) and the joint KMS part of (Q) follow from one premise with no temperature in it, posable only on the continuum}(iii) the generators' modes are then jointly KMS for the Killing flow at $\beta = 8\pi M$. The wording above, on the ladders $\{v_{k}, v_{k}^{*}\}$ with a phase-acting flow, is a discretisation of that statement, recovered only for modes of vanishing Mellin width; (N$_{+}$) itself cannot be posed on point spectrum. A hole formed by collapse is not covered by (N$_{+}$), and (R) remains a premise. (The whole generator of the eternal horizon, which that case of (N$_{+}$) uses, lies beyond one flat chart; see that premise.)
\end{postulate}

\section{Setting and imports}

\phantomsection\label{c:note.hawking.setting-setting-imports}
Units are $G = c = \hbar = k_B = 1$ throughout, so for Schwarzschild the surface gravity is $\kappa = 1/4M$, the Hawking temperature is $\TH = \kappa/2\pi = 1/8\pi M$, the modular period is $\beta := 1/\TH = 8\pi M$, and areas are in Planck units, $\ell_P = 1$. (Two symbols are nonetheless kept visible where they carry the point: $\hbar$ in the section ``The four-dimensional lift'', where the ordering shift is a multiple of it, and $\ell_P$ in the section ``The finite boundary mode space, and the $\ell_P^2/A$ correction'', where the corrections scale as $\ell_P^2/A$, $A$ the horizon area.)

\subsection*{The single mode, self-contained}
\phantomsection\label{c:note.hawking.setting-single-mode-self-contained}
The unit of the boundary mode space is a completed null pair: vectors $v, v^*$ with $v^2 = (v^*)^2 = 0$ and $vv^* + v^*v = 1$, the ladder pair the transport principle adjoins \cite{ButlerT1}. Everything modular about a single mode is elementary and is proved here outright, per the series' standard.

\begin{proposition}[Modular structure of the completed pair]\label{c:hawking.modular-structure-completed-pair}
\changedmark{gray}{reworded}
Let $(v, v^*)$ be a completed pair and $\sigma := v^*v - vv^*$. Then:
\begin{enumerate}
\item $\sigma^2 = 1$; in the hyperbolic atom $\Cl(1,1)$ with $v = \tfrac12(e_0 + e_1)$, $v^* = \tfrac12(e_0 - e_1)$, one has $\sigma = e_0 e_1$: the unit boost bivector of the plane (relabelling the pair reverses the orientation);
\item the rotor flow $R_\lambda := \exp(\lambda\sigma/2)$ acts by $R_\lambda\, v\, R_\lambda^{-1} = e^{-\lambda}v$ and $R_\lambda\, v^*\, R_\lambda^{-1} = e^{+\lambda}v^*$: the pair diagonalises the boost with opposite weights;
\item the spectral projectors $\tfrac12(1 \pm \sigma)$ are the occupation projectors $v^*v$ and $vv^*$;
\item in the Fock representation the boost generator is the shifted number operator, $\tfrac12\sigma = N - \tfrac12$;
\item the bivector space of $\Cl(1,1)$ is one-dimensional, so the boost flow is the unique one-parameter rotor flow of the atom \emph{up to reparametrisation of the flow parameter and orientation}: the algebra fixes the orbits of any rotor-implemented modular flow, but not its speed, which is exactly the normalisation Remark \ref{c:note.hawking.what-now-proved-open-problem} (the section ``The state input: the I-state'') leaves open.
\end{enumerate}
\end{proposition}
\begin{proof}
(1) $\sigma^2 = v^*vv^*v + vv^*vv^* - v^*v\,vv^* - vv^*\,v^*v = v^*v + vv^* - 0 - 0 = 1$, using $vv^*v = v$, $v^*vv^* = v^*$ and the vanishing squares. In $\Cl(1,1)$, by \ref{c:std-algebra.vector-squares}: $v^*v = \tfrac14(e_0 - e_1)(e_0 + e_1) = \tfrac14(e_0^2 - e_1^2 + e_0e_1 - e_1e_0) = \tfrac14(1 + 1 + 2e_0e_1) = \tfrac12(1 + e_0e_1)$, so $\sigma = v^*v - vv^* = e_0e_1$ in this orientation; relabelling the pair reverses it. (2) From $vv^*v = v$ and the vanishing squares, $\sigma v = (v^*v - vv^*)v = -vv^*v = -v$ and $v\sigma = v(v^*v - vv^*) = vv^*v = +v$. The first relation gives $e^{\lambda\sigma/2}\,v = v\,e^{-\lambda\sigma/2}$ term by term in the exponential series, so $R_\lambda\, v\, R_\lambda^{-1} = v\,e^{-\lambda\sigma}$. Since $\sigma^2 = 1$, $e^{-\lambda\sigma} = \cosh\lambda - \sigma\sinh\lambda$ (\ref{c:std-algebra.bivector-exponential}), and the second relation collapses it on the left factor: $v\,e^{-\lambda\sigma} = v(\cosh\lambda - \sinh\lambda) = e^{-\lambda}v$. The $v^*$ computation is the mirror image, with both signs reversed. (3) $\tfrac12(1 + \sigma) = \tfrac12(1 + v^*v - vv^*) = \tfrac12(v^*v + vv^* + v^*v - vv^*) = v^*v$, and likewise for the complement. (4) In the two-dimensional Fock module ($v^* = $ creation, $v = $ annihilation, $N = v^*v$): $\tfrac12\sigma = \tfrac12(v^*v - vv^*) = \tfrac12(N - (1 - N)) = N - \tfrac12$. (5) The bivector space of $\Cl(1,1)$ is the line $\R\,e_0e_1$, so every one-parameter rotor group is of the form $\lambda \mapsto \exp(c\lambda\sigma/2)$ for some real $c$: the same orbits at a rescaled parameter, which is the reparametrisation the statement reserves. The verification script checks every clause in exact arithmetic, including the finite conjugation of clause (2) at a rational rapidity.

\end{proof}

\begin{fact}[Imported: modular descent, its premises, and the geometric normalisation (external ingredients)]\label{c:hawking.imported-modular-descent-premises-geometric-normalisation-external}
The results from outside the series on which \ref{c:hawking.imported-modular-descent-premises-geometric-normalisation} draws, as cited there: J. J. Bisognano and E. H. Wichmann, On the duality condition for quantum fields, J. Math. Phys. 17 (1976), 303--321; G. L. Sewell, Quantum fields on manifolds: PCT and gravitationally induced thermal states, Ann. Phys. 141 (1982), 201--224; W. Israel, Thermo-field dynamics of black holes, Phys. Lett. A 57 (1976), 107--110.
\emph{Source:} J. J. Bisognano and E. H. Wichmann, On the duality condition for quantum fields, J. Math. Phys. 17 (1976), 303--321.; G. L. Sewell, Quantum fields on manifolds: PCT and gravitationally induced thermal states, Ann. Phys. 141 (1982), 201--224.; W. Israel, Thermo-field dynamics of black holes, Phys. Lett. A 57 (1976), 107--110..
\end{fact}

\begin{assembled}[Imported: modular descent, its premises, and the geometric normalisation]\label{c:hawking.imported-modular-descent-premises-geometric-normalisation}
\changedmark{BurntOrange}{changed since the deposit}
Let the ascended algebra of a horizon pair carry a state that is pure, invariant under the boost flow of Proposition \ref{c:hawking.modular-structure-completed-pair}, and entangling across the pair (cyclic and separating for the one-sided subalgebra). Then the descended state is KMS for the phase flow $\operatorname{Ad}(e^{it\sigma/2})$ of the same generator at \emph{some} inverse temperature $\beta \in \R$ of rapidity, and the hypotheses leave $\beta$ free, its sign included; the Gibbs \emph{form} of the descended state is state-independent. Positivity and the value are supplied by T5's named premises: under (P) univalence and (R) the KMS condition at the geometric period (one modular period equal to one rapidity period $2\pi$, i.e.\ Schwarzschild time $2\pi/\kappa$, the Euclidean period $8\pi M$), the descended state is Fermi-Dirac at $\beta = 2\pi$ of rapidity, which is the physical temperature $\TH = \kappa/2\pi = 1/8\pi M$; the boost plane is the one the shape operator scales by $-2$, and $\kappa = 1/4M$ is proved from the gauge field. Per mode of Schwarzschild frequency $\omega$ the reduced density matrix is the Fermi-Dirac matrix $\mathrm{diag}(1, e^{-\omega/\TH})/Z$. \emph{Provenance:} T5 \cite[Thms.~7.3 and 7.6, Prems.~7.8 and 7.19, Prop.~7.18, Thm.~7.20, Cor.~6.4]{ButlerT5}, with $\kappa = 1/4M$ from the gauge field in \cite[App.~A]{ButlerT5}; premise (R) is proved irreducible there, the temperature being no algebra-intrinsic quantity \cite[Thm.~7.16]{ButlerT5}. The wedge-modular mechanism is Bisognano-Wichmann-Sewell \cite{BisognanoWichmann1976, Sewell1982}, and the thermofield reading is Israel's \cite{Israel1976}.
\end{assembled}
\begin{proof}
Let the ascended algebra of a horizon pair carry a state that is pure, invariant under the boost flow of Proposition \ref{c:hawking.modular-structure-completed-pair}, and entangling across the pair (cyclic and separating for the one-sided subalgebra). Then the descended state is KMS for the phase flow $\operatorname{Ad}(e^{it\sigma/2})$ of the same generator at \emph{some} inverse temperature $\beta \in \R$ of rapidity, and the hypotheses leave $\beta$ free, its sign included; the Gibbs \emph{form} of the descended state is state-independent. Positivity and the value are supplied by T5's named premises: under (P, \ref{c:singularity.p-univalence}) univalence and (R, \ref{c:singularity.r-regularity-fixed-locus}) the KMS condition at the geometric period (one modular period equal to one rapidity period $2\pi$, i.e.\ Schwarzschild time $2\pi/\kappa$, the Euclidean period $8\pi M$), the descended state is Fermi-Dirac at $\beta = 2\pi$ of rapidity, which is the physical temperature $\TH = \kappa/2\pi = 1/8\pi M$; the boost plane is the one the shape operator scales by $-2$, and $\kappa = 1/4M$ is proved from the gauge field. Per mode of Schwarzschild frequency $\omega$ the reduced density matrix is the Fermi-Dirac matrix $\mathrm{diag}(1, e^{-\omega/\TH})/Z$. \emph{Provenance:} T5 \cite[Thms.~7.3 and 7.6, Prems.~7.8 and 7.19, Prop.~7.18, Thm.~7.20, Cor.~6.4]{ButlerT5} (\ref{c:singularity.modular-descent}, \ref{c:singularity.euclidean-section-period}, \ref{c:singularity.rate-discharge-hypothesis-iii}, \ref{c:singularity.matching-steps-4b-4c-join-equation}, \ref{c:singularity.descended-form-state-independent}), with $\kappa = 1/4M$ from the gauge field in \cite[App.~A]{ButlerT5}; premise (R) is proved irreducible there, the temperature being no algebra-intrinsic quantity \cite[Thm.~7.16]{ButlerT5} (\ref{c:singularity.rate-rigidity-temperature-not-algebra-intrinsic}). The wedge-modular mechanism is Bisognano-Wichmann-Sewell \cite{BisognanoWichmann1976, Sewell1982}, and the thermofield reading is Israel's \cite{Israel1976}.

This restates \ref{c:hawking.imported-modular-descent-premises-geometric-normalisation-external}, \ref{c:singularity.computed}.

\end{proof}

\begin{assembled}[Imported: the algebra of the boundary modes is the CAR algebra]\label{c:hawking.imported-algebra-boundary-modes-car-algebra}
\changedmark{gray}{reworded}
For even transport charge, the completed complexified limit of the tower of iterated completions is $\Cl(p,q)_\C \otimes \CAR$, the ground algebra tensored with the CAR algebra of countably many fermionic modes, and the mode operators are the adjoined null pairs themselves ($a_j = v_j^* = w_j$, $a_j^\dagger = v_j$, with the Fock adjointness). Here the \emph{transport charge} is the mod-8 invariant the tower's completions advance and the \emph{lineage} is the tower of iterated completions itself, both as in \cite{ButlerT3}; the physical lineage has charge six, so the theorem applies. \emph{Provenance:} the tower paper \cite[Prop.~4.1, Thms.~5.3 and 5.6]{ButlerT3}.
\end{assembled}
\begin{proof}
For even transport charge, the completed complexified limit of the tower of iterated completions is $\Cl(p,q)_\C \otimes \CAR$, the ground algebra tensored with the CAR algebra of countably many fermionic modes, and the mode operators are the adjoined null pairs themselves ($a_j = v_j$, $a_j^\dagger = v_j^*$, with the Fock adjointness). Here the \emph{transport charge} is the mod-8 invariant the tower's completions advance and the \emph{lineage} is the tower of iterated completions itself, both as in \cite{ButlerT3}; the physical lineage has charge six, so the theorem applies. \emph{Provenance:} the tower paper \cite[Prop.~4.1, Thms.~5.3 and 5.6]{ButlerT3} (\ref{c:tower.algebraic-limit-even-charge}, \ref{c:tower.car-limit-even-charge}, \ref{c:tower.ladder-pairs-satisfy-car-vectors}).

\end{proof}

\begin{fact}[Imported: the s-wave bosonic sector (external ingredients)]\label{c:hawking.imported-s-wave-bosonic-sector-external}
The results from outside the series on which \ref{c:hawking.imported-s-wave-bosonic-sector} draws, as cited there: I. B. Frenkel and V. G. Kac, Basic representations of affine Lie algebras and dual resonance models, Invent. Math. 62 (1980), 23--66.
\emph{Source:} I. B. Frenkel and V. G. Kac, Basic representations of affine Lie algebras and dual resonance models, Invent. Math. 62 (1980), 23--66..
\end{fact}

\begin{assembled}[Imported: the s-wave bosonic sector]\label{c:hawking.imported-s-wave-bosonic-sector}
\changedmark{gray}{reworded}
Two halves, at two levels. \emph{Unconditional:} on a window of the boundary mode space, the number-conserving bilinears (currents) satisfy the free-boson commutation relations exactly in the completed limit, with an exact finite-window operator identity whose deviation is a difference of edge operators; and the bilinear subalgebra is globally invariant under the modular boost. \emph{Conditional on T5's Premise (Q), joint quasi-freeness with the s-wave weights:} the restriction of the Fermi-Dirac KMS state to the bilinear subalgebra is KMS at the same temperature (an instance of Lemma \ref{c:hawking.kms-restricts-invariant-subalgebras} below), giving the Bose-Einstein occupancy at $\TH$ for the s-wave sector; the linear weight assignment is the s-wave reduction's convention, adopted there as part of the premise. \emph{Provenance:} T5 \cite[Prem.~8.1, Thm.~8.2]{ButlerT5}; the correspondence itself is classical \cite{FrenkelKac1980}. The exact edge identity is reproved, in the form this paper needs, as Proposition \ref{c:hawking.current-commutator-exactly-edge-split} below.
\end{assembled}
\begin{proof}
Two halves, at two levels. \emph{Unconditional:} on a window of the boundary mode space, the number-conserving bilinears (currents) satisfy the free-boson commutation relations exactly in the completed limit, with an exact finite-window operator identity whose deviation is a difference of edge operators; and the bilinear subalgebra is globally invariant under the modular boost. \emph{Conditional on T5's Premise (Q) (\ref{c:singularity.q-joint-quasi-freeness-s-wave}), joint quasi-freeness with the s-wave weights:} the restriction of the Fermi-Dirac KMS state to the bilinear subalgebra is KMS at the same temperature (an instance of Lemma \ref{c:hawking.kms-restricts-invariant-subalgebras} below), giving the Bose-Einstein occupancy at $\TH$ for the s-wave sector; the linear weight assignment is the s-wave reduction's convention, adopted there as part of the premise. \emph{Provenance:} T5 \cite[Prem.~8.1, Thm.~8.2]{ButlerT5} (\ref{c:singularity.bosonisation-bridge-s-wave-sector}); the correspondence itself is classical \cite{FrenkelKac1980}. The exact edge identity is reproved, in the form this paper needs, as Proposition \ref{c:hawking.current-commutator-exactly-edge-split} below.

This restates \ref{c:hawking.imported-s-wave-bosonic-sector-external}.

\end{proof}

\begin{assembled}[Imported: the exterior is exactly general relativity]\label{c:hawking.imported-exterior-exactly-general-relativity}
On the region where the gauge field is invertible, and at $\Lambda = 0$, the series' description coincides exactly with classical general relativity: the field equations are equivalent to the Einstein equations, and the vacuum equations are exactly the covariant constancy of the scale tractor. \emph{Provenance:} \cite[Fact 2.1, Thms.~4.5 and 4.7]{ButlerT6} (Thm.~4.7 is the complementary $\Lambda \neq 0$ case), resting on Lasenby-Doran-Gull's equivalence theorem.
\end{assembled}
\begin{proof}
On the region where the gauge field is invertible, and at $\Lambda = 0$, the series' description coincides exactly with classical general relativity: the field equations are equivalent to the Einstein equations, and the vacuum equations are exactly the covariant constancy of the scale tractor. \emph{Provenance:} \cite[Fact 2.1, Thms.~4.5 and 4.7]{ButlerT6} (\ref{c:interior.covariant-form-vacuum-equations}, \ref{c:interior.imported-equivalence-einstein-equations}) (Thm.~4.7 is the complementary $\Lambda \neq 0$ case), resting on Lasenby-Doran-Gull's equivalence theorem.

\end{proof}

\begin{assembled}[Imported: the conformal storey (twistor storey), and the ordering shift]\label{c:hawking.imported-twistor-storey-ordering-shift}
\changedmark{gray}{reworded}
The spinor module of the completed algebra $\Cl(2,4)$ is the conformal spinor space $\CS = \Cl^+(2,4)P$ (twistor space), of dimension $4$ over the complex unit $I_6$ (a \emph{storey} is the series' word for one level of the completion tower, as in \cite{ButlerT9}), the palatial commutation relation $[Z^\mu, \bar Z_\nu] = \hbar\,\delta^\mu_\nu$ on the quadratic (palatial) quantisation of the storey ($\bar Z$ the dual of $Z$ under the reversion form $h$), and, for $M \in \mathfrak{gl}(\CS)$, the symmetrised quadratic $q(M) := \tfrac12\bigl(Z^\mu M_\mu^{\ \nu}\bar Z_\nu + \bar Z_\nu M_\mu^{\ \nu} Z^\mu\bigr)$. (Indices on $\CS$ are written $\mu, \nu$ throughout; $\alpha$ and $\beta$ are reserved for the modular dynamics and its inverse temperature.) Symmetric ordering shifts the distinguished $\mathfrak{u}(1)$ generator $q(\id)$ by $-\hbar/2$ per complex dimension of $\CS$, totalling $-2\hbar$ (equivalently, $q(M) = \sum M_\mu^{\ \nu} Z^\mu \bar Z_\nu - \tfrac{\hbar}{2}\tr M$); on the Clifford storey, the boost generator of each adjoined plane is $N - \tfrac12$, a shift of minus one half per fermionic mode: the two storeys carry the same signed ordering shift. \emph{Provenance:} the correspondence paper \cite[Thm.~3.1, Rem.~3.2, \S 5, Lem.~5.5, Prop.~5.6, Thm.~5.7]{ButlerT9}; the homogeneity consequence is computed inside the proof of its Theorem 5.7 and is promoted to a stated result as Proposition \ref{c:hawking.homogeneity-ordering-shift} below.
\end{assembled}
\begin{proof}
The spinor module of the completed algebra $\Cl(2,4)$ is the conformal spinor space $\CS = \Cl^+(2,4)P$ (twistor space), of dimension $4$ over the complex unit $I_6$ (a \emph{storey} is the series' word for one level of the completion tower, as in \cite{ButlerT9}), the palatial commutation relation $[Z^\mu, \bar Z_\nu] = \hbar\,\delta^\mu_\nu$ on the quadratic (palatial) quantisation of the storey ($\bar Z$ the dual of $Z$ under the reversion form $h$), and, for $M \in \mathfrak{gl}(\CS)$, the symmetrised quadratic $q(M) := \tfrac12\bigl(Z^\mu M_\mu^{\ \nu}\bar Z_\nu + \bar Z_\nu M_\mu^{\ \nu} Z^\mu\bigr)$. (Indices on $\CS$ are written $\mu, \nu$ throughout; $\alpha$ and $\beta$ are reserved for the modular dynamics and its inverse temperature.) Symmetric ordering shifts the distinguished $\mathfrak{u}(1)$ generator $q(\id)$ by $-\hbar/2$ per complex dimension of $\CS$, totalling $-2\hbar$ (equivalently, $q(M) = \sum M_\mu^{\ \nu} Z^\mu \bar Z_\nu - \tfrac{\hbar}{2}\tr M$); on the Clifford storey, the boost generator of each adjoined plane is $N - \tfrac12$, a shift of minus one half per fermionic mode: the two storeys carry the same signed ordering shift. \emph{Provenance:} the correspondence paper \cite[Thm.~3.1, Rem.~3.2, \S 5, Lem.~5.5, Prop.~5.6, Thm.~5.7]{ButlerT9} (\ref{c:correspondence.fock-model-complexified-ascended-algebra}, \ref{c:correspondence.clifford-side-ordering-shift}, \ref{c:correspondence.symmetrised-quadratics-palatial-algebra}, \ref{c:correspondence.shared-core-ordering-dictionary}); the homogeneity consequence is computed inside the proof of its Theorem 5.7 and is promoted to a stated result as Proposition \ref{c:hawking.homogeneity-ordering-shift} below.

\end{proof}

\begin{definition}[The boundary-mode algebra, windows, and the boost dynamics]\label{c:obj.hawking.boundary-mode-algebra-windows-boost-dynamics}
\changedmark{BurntOrange}{changed since the deposit}
Fix a countable mode set $J$ with Schwarzschild frequencies $\omega_j > 0$ ($j \in J$), the frequencies the geometry assigns to the boundary modes (Fact \ref{c:hawking.imported-modular-descent-premises-geometric-normalisation}). \emph{Provenance of $J$:} the boundary modes are the ladder pairs of the blown-up singular boundary, one per point, aggregated over the s-wave channels into the countable array of Theorem \ref{c:singularity.bosonisation-bridge-s-wave-sector}, an identification the singularity paper makes and flags. By convention here (it is in neither the singularity paper nor the thermality paper, and reconciles the array's weights $\kappa k$, $k \in \mathbb{Z}$, with $\omega_j > 0$), the modes below the level enter $J$ through their particle-hole conjugates, so that every frequency is positive, and the one zero-frequency mode of the array lies outside $J$; the bookkeeping is filed alongside (synthesis.s-wave-array-is-a-boundary-mode-set). The tower of iterated completions is the algebra's model, not the source of its countability: its limit is the CAR algebra with the adjoined null pairs as modes (Fact \ref{c:hawking.imported-algebra-boundary-modes-car-algebra}), while on the physical closure one ascent completes the description and further rungs lie on shadow walls outside the physical region (Corollary \emph{Answers to the iteration question} [T3, Cor 6.2]). The \emph{boundary-mode algebra} $\mathfrak{A}$ is the CAR algebra over $\ell^2(J)$ (Fact \ref{c:hawking.imported-algebra-boundary-modes-car-algebra}), the C*-inductive limit of the \emph{window algebras} $\mathfrak{A}_F$, $F \subset J$ finite, each $\mathfrak{A}_F \cong \End(\Lambda^\bullet \C^F)$ of dimension $4^{|F|}$. The \emph{boost dynamics} is the one-parameter group $\tau_t$ acting on generators by $\tau_t(a_j) = e^{-i\omega_j t}\,a_j$, with $t$ Schwarzschild time: the quasi-free flow whose one-particle Hamiltonian is $h = \mathrm{diag}(\omega_j)$; it preserves every window. In rapidity $\lambda = \kappa t$ the same flow has boost weights $\omega_j/\kappa$ and modular temperature $2\pi$ (Fact \ref{c:hawking.imported-modular-descent-premises-geometric-normalisation}); in Schwarzschild time the inverse temperature is $\beta = 1/\TH = 2\pi/\kappa$, and this normalisation is used throughout.
\end{definition}

\section{The state input: the I-state}

\phantomsection\label{c:note.hawking.setting-state-input-i-state}
T5 proves single-mode thermality from its premises (P) and (R), deriving the occupancy from the gluing identity, and says exactly what is premised. This section lifts the statement to the full mode algebra: within the quasi-free class the state is unique and its two-point data are forced outright, all modes at once; given the dynamics, even the class restriction can be removed, by a theorem of the literature; what cannot yet be derived is the dynamics-selection axiom itself, and that open problem is posed at the end, against T5's own selection analysis.

\subsection*{Two elementary lemmas}
\phantomsection\label{c:note.hawking.setting-two-elementary-lemmas}
We write $\alpha_t$ for a finite-dimensional dynamics; on a window $F$ the relevant instance is $\alpha_t = \tau_t|_{\mathfrak{A}_F}$, generated by $H_F$, the second quantisation of $h|_F$.

\begin{lemma}[Gibbs states are KMS, and conversely, in finite dimensions]\label{c:hawking.gibbs-states-are-kms-conversely-finite}
Let $\mathfrak{B} = \End(\C^m)$, let $H$ be diagonal in the basis $(\xi_a)$ with eigenvalues $\lambda_a$, all distinct or not, let $\alpha_t = \mathrm{Ad}(e^{itH})$, and $\beta \in \R$. Then the Gibbs state $\omega_G(x) := \tr(e^{-\beta H}x)/\tr(e^{-\beta H})$ satisfies the KMS identity
\[
\omega_G(x\,y) \;=\; \omega_G\bigl(y\,\alpha_{i\beta}(x)\bigr), \qquad \alpha_{i\beta}(x) := e^{-\beta H}x\,e^{\beta H},
\]
for all $x, y$; and it is the \emph{unique} $\alpha$-KMS state at $\beta$ on $\mathfrak{B}$, degenerate spectrum included.
\end{lemma}
\begin{proof}
The identity is cyclicity of the trace:
\[
\omega_G\bigl(y\,\alpha_{i\beta}(x)\bigr) = \tr\bigl(e^{-\beta H}y\,e^{-\beta H}x\,e^{\beta H}\bigr)/Z = \tr\bigl(e^{-\beta H}xy\bigr)/Z = \omega_G(xy),
\]
the trailing $e^{\beta H}$ moved around the trace to cancel the inner factor. Conversely write $\omega(x) = \tr(\rho x)$ and evaluate the KMS identity on matrix units $x = E_{ab}$, $y = E_{cd}$: since $\alpha_{i\beta}(E_{ab}) = e^{-\beta(\lambda_a - \lambda_b)}E_{ab}$ and $\tr(\rho\,E_{ab}E_{cd}) = \delta_{bc}\,\rho_{da}$, the identity reads
\[
\delta_{bc}\,\rho_{da} \;=\; e^{-\beta(\lambda_a - \lambda_b)}\,\delta_{da}\,\rho_{bc}
\qquad \text{for all } a, b, c, d.
\]
Taking $c = b$ and $d = a$ gives $\rho_{aa} = e^{-\beta(\lambda_a - \lambda_b)}\rho_{bb}$: Gibbs ratios between all diagonal entries. Taking $c = b$ and $d \neq a$ gives $\rho_{da} = 0$ (the right side carries $\delta_{da} = 0$), for every pair $d \neq a$, degenerate eigenvalues included: $\rho$ is diagonal. So $\rho = e^{-\beta H}/Z$, and the KMS state is unique. The verification script performs the solve symbolically for the mode algebra $\Mat_2$ with $H = \omega N$ and confirms both clauses as identities of rational functions in the Boltzmann factor.

\end{proof}

\begin{lemma}[KMS restricts to invariant subalgebras]\label{c:hawking.kms-restricts-invariant-subalgebras}
If $\omega$ is $\alpha$-KMS at $\beta$ on $\mathfrak{B}$ and $\mathfrak{C} \subseteq \mathfrak{B}$ is a unital subalgebra with $\alpha_t(\mathfrak{C}) = \mathfrak{C}$, then $\omega|_{\mathfrak{C}}$ is KMS at $\beta$ for $\alpha|_{\mathfrak{C}}$.
\end{lemma}
\begin{proof}
For $x, y \in \mathfrak{C}$ the correlation function $t \mapsto \omega(y\,\alpha_t(x))$ is computed entirely from elements of $\mathfrak{C}$ (invariance keeps $\alpha_t(x) \in \mathfrak{C}$), and is the same function of $t$ whether $x, y$ are regarded in $\mathfrak{C}$ or in $\mathfrak{B}$. Its continuation to the closed strip is unique (a function analytic on the strip and continuous on its closure is determined by its boundary restriction to the real line, by uniqueness of analytic continuation), so the strip function and the KMS boundary identity transfer verbatim from $\mathfrak{B}$ to $\mathfrak{C}$.

\end{proof}

\begin{fact}[Imported: quasi-free states of the CAR algebra]\label{c:hawking.imported-quasi-free-states-car-algebra}
For every operator $A$ on $\ell^2(J)$ with $0 \le A \le 1$ (the \emph{symbol}) there is a unique gauge-invariant state $\omega_A$ on $\mathfrak{A}$ with two-point function $\omega_A(a^\dagger(f)\,a(g)) = \ip{g}{Af}$ and all higher truncated functions zero (Wick determinants). More generally, a quasi-free state (all truncated functions beyond second order vanishing) is determined by its full two-point data, the symbol together with the pairing $P(f, g) := \omega(a^\dagger(f)\,a^\dagger(g))$, and gauge invariance is exactly $P = 0$. On a finite window with diagonal symbol this is the explicit product state $\bigotimes_j \rho_j$, $\rho_j = \mathrm{diag}(1 - n_j,\, n_j)$ in the occupation basis, which the script realises exactly. \emph{Provenance:} Araki \cite{Araki1971}; \cite[\S 5.2]{BratteliRobinson1997}.
\emph{Source:} H. Araki, On quasifree states of CAR and Bogoliubov automorphisms, Publ. RIMS Kyoto Univ. 6 (1971), 385--442.; O. Bratteli and D. W. Robinson, Operator Algebras and Quantum Statistical Mechanics 2, 2nd ed., Springer, Berlin (1997)., \S 5.2.
\end{fact}

\begin{fact}[Imported: KMS passes to inductive limits]\label{c:hawking.imported-kms-passes-inductive-limits}
Let a C*-algebra be the closure of an increasing union of subalgebras, each globally invariant under a one-parameter dynamics, and let a state be KMS at $\beta$ on each subalgebra for the restricted dynamics. Then the state is KMS at $\beta$ on the whole algebra: the analytic-element requirement is checked on the dense union and passes to the closure. \emph{Provenance:} \cite[Prop.~5.3.25]{BratteliRobinson1997}.
\emph{Source:} O. Bratteli and D. W. Robinson, Operator Algebras and Quantum Statistical Mechanics 2, 2nd ed., Springer, Berlin (1997)., Prop.~5.3.25.
\end{fact}

\begin{fact}[Imported: quasi-free KMS reduces to the two-point level]\label{c:hawking.imported-quasi-free-kms-reduces-two}
For the CAR algebra over a one-particle space with the quasi-free dynamics of a self-adjoint one-particle Hamiltonian, a gauge-invariant quasi-free state is KMS at $\beta$ if and only if its two-point function satisfies the two-point KMS identity: through the Wick-determinant structure, the analytic-continuation requirement for products of generators reduces to the two-point level. \emph{Provenance:} Araki \cite{Araki1971}; \cite[\S 5.3]{BratteliRobinson1997}.
\emph{Source:} H. Araki, On quasifree states of CAR and Bogoliubov automorphisms, Publ. RIMS Kyoto Univ. 6 (1971), 385--442.; O. Bratteli and D. W. Robinson, Operator Algebras and Quantum Statistical Mechanics 2, 2nd ed., Springer, Berlin (1997)., \S 5.3.
\end{fact}

\begin{theorem}[The I-state]\label{c:hawking.i-state}
Consider the boundary-mode algebra with its boost dynamics (Definition \ref{c:obj.hawking.boundary-mode-algebra-windows-boost-dynamics}) at inverse temperature $\beta = 1/\TH = 2\pi/\kappa$.
\begin{enumerate}
\item \emph{(Existence, exactly on windows.)} The quasi-free state $\omega_I$ with the Fermi-Dirac symbol
\[
A_I \;=\; \bigl(1 + e^{\beta h}\bigr)^{-1}, \qquad\text{i.e.}\qquad n_j = \frac{1}{e^{\omega_j/\TH} + 1},
\]
restricts on every window to the Gibbs state of $H_F = \sum_{j \in F}\omega_j N_j$ and is KMS there, exactly (Lemma \ref{c:hawking.gibbs-states-are-kms-conversely-finite}); it is the unique quasi-free extension of these window states, and it is $\tau$-KMS at $\beta$ on the full algebra by Fact \ref{c:hawking.imported-kms-passes-inductive-limits}.
\item \emph{(Gauge invariance is forced.)} If a quasi-free state of $\mathfrak{A}$ is $\tau$-KMS at $\beta$, its pairing vanishes: the state is gauge-invariant automatically.
\item \emph{(The symbol is derived.)} Its symbol then satisfies the operator identity $A(1 + e^{-\beta h}) = e^{-\beta h}$, equivalently $A = e^{-\beta h}(1 - A)$, whence $A = A_I$: within the quasi-free class, the state whose modular clock is the geometry's clock is unique, and its two-point data are forced.
\end{enumerate}
\end{theorem}
\begin{proof}
(1) On a window $F$, the quasi-free state with diagonal symbol is the product state of Fact \ref{c:hawking.imported-quasi-free-states-car-algebra} with $n_j = (1 + e^{\beta\omega_j})^{-1}$, which is exactly the Gibbs density matrix of $H_F$: $e^{-\beta H_F} = \bigotimes_j \mathrm{diag}(1, e^{-\beta\omega_j})$, normalised; the script builds the two sides independently and confirms the identity in exact arithmetic. Lemma \ref{c:hawking.gibbs-states-are-kms-conversely-finite} gives the KMS identity on the window, exactly; the script verifies it on a three-mode window at rational Boltzmann factors, on a spanning set of pairs. The windows are nested with dense union and the dynamics preserves each, so Fact \ref{c:hawking.imported-kms-passes-inductive-limits} passes the KMS property to $\mathfrak{A}$.
(2) Apply the KMS identity to the pair $x = a^\dagger(f)$, $y = a^\dagger(g)$, with $f, g$ supported in windows (so that $\alpha_{i\beta}(a^\dagger(f)) = a^\dagger(e^{-\beta h}f)$ is again a generator):
\[
P(f, g) \;=\; \omega\bigl(a^\dagger(f)\,a^\dagger(g)\bigr) \;=\; \omega\bigl(a^\dagger(g)\,a^\dagger(e^{-\beta h}f)\bigr) \;=\; P(g,\, e^{-\beta h}f).
\]
Antisymmetry of creation operators gives $P(f, g) = -P(g, f)$, so $P(g, (1 + e^{-\beta h})f) = 0$ for all $f, g$ in the dense span of the windows; $1 + e^{-\beta h}$ is bounded with bounded inverse and preserves that span, so $P = 0$.
(3) Apply the KMS identity to the pair $x = a^\dagger(f)$, $y = a(g)$:
\[
\omega\bigl(a^\dagger(f)\,a(g)\bigr) \;=\; \omega\bigl(a(g)\,a^\dagger(e^{-\beta h}f)\bigr)
\;=\; \ip{g}{e^{-\beta h}f} - \omega\bigl(a^\dagger(e^{-\beta h}f)\,a(g)\bigr),
\]
using the CAR $\{a(g), a^\dagger(f')\} = \ip{g}{f'}$. In symbols: $\ip{g}{Af} = \ip{g}{e^{-\beta h}f} - \ip{g}{Ae^{-\beta h}f}$ for all $f, g$ supported in windows, i.e.\ $A(1 + e^{-\beta h}) = e^{-\beta h}$ on the dense span, i.e.\ $A = e^{-\beta h}(1+e^{-\beta h})^{-1} = (1 + e^{\beta h})^{-1} = A_I$ (all four operators bounded, $h \ge 0$ here). By (2) and Fact \ref{c:hawking.imported-quasi-free-states-car-algebra}, the state is determined by its two-point data, which are forced, so $\omega = \omega_I$. The mode-diagonal scalar identity $n = e^{-\beta\omega}(1-n)$, equivalently the detailed balance $n/(1-n) = e^{-\beta\omega}$, is checked by the script as an exact identity of rational functions.

\end{proof}

\begin{fact}[Imported: ideal-Fermi-gas uniqueness]\label{c:hawking.imported-ideal-fermi-gas-uniqueness}
For the CAR algebra over a one-particle space with self-adjoint $h$ and the induced quasi-free dynamics, the $\beta$-KMS state is unique and equals the quasi-free state with symbol $(1 + e^{\beta h})^{-1}$. \emph{Provenance:} \cite[Ex.~5.3.2 and \S 5.3]{BratteliRobinson1997}; the quasi-free calculus is \cite{Araki1971}.
\emph{Source:} O. Bratteli and D. W. Robinson, Operator Algebras and Quantum Statistical Mechanics 2, 2nd ed., Springer, Berlin (1997)., Ex.~5.3.2 and \S 5.3; H. Araki, On quasifree states of CAR and Bogoliubov automorphisms, Publ. RIMS Kyoto Univ. 6 (1971), 385--442..
\end{fact}

\begin{corollary}[Full uniqueness, given the dynamics]\label{c:hawking.full-uniqueness-given-dynamics}
Among all states of $\mathfrak{A}$, $\omega_I$ is the unique $\tau$-KMS state at $\beta$.
\end{corollary}
\begin{proof}
Fact \ref{c:hawking.imported-ideal-fermi-gas-uniqueness}, whose hypotheses Definition \ref{c:obj.hawking.boundary-mode-algebra-windows-boost-dynamics} satisfies; the quasi-free state it names has symbol $A_I$, which is $\omega_I$ by Fact \ref{c:hawking.imported-quasi-free-states-car-algebra}.

\end{proof}

\begin{remark}[What is now proved, and the open problem]\label{c:note.hawking.what-now-proved-open-problem}
Three levels, stated separately. \emph{Proved here:} within the quasi-free class the state is unique, with gauge invariance and the Fermi-Dirac symbol both forced by the KMS identity (Theorem \ref{c:hawking.i-state}); in this sense the I-state is the framework's equilibrium vacuum, and the quasi-free restriction is a convenience of proof, removed at the next level. \emph{Imported:} given the boost dynamics and $\beta$, full uniqueness among all states (Corollary \ref{c:hawking.full-uniqueness-given-dynamics}, from Fact \ref{c:hawking.imported-ideal-fermi-gas-uniqueness}). \emph{Open, and named:} deriving the dynamics itself, i.e.\ an axiom selecting the boost flow and its temperature from the geometry. T5 has already treated this question at the level of a single factor: its geometric modular condition (G) is formulated and rejected there as over-determining the boundary modes, its premise (R) is proved irreducible, its rate-rigidity theorem shows the temperature is not algebra-intrinsic, and the resulting state class is recorded \cite[Def.~7.12, Thms.~7.16 and 7.28, Cor.~7.32]{ButlerT5}. The question this paper poses is different and remains open: a \emph{net-level} selection in the style of geometric modular action \cite{BDFS2000}, demanding that the modular groups of a net of boundary-mode algebras act geometrically and deriving the boost identification for the whole net rather than assuming it. Within the single atom the algebra already forces the orbits of any rotor-implemented flow, though not its speed (Proposition \ref{c:hawking.modular-structure-completed-pair}(5)); the execution for the net is a successor problem (the section ``Scope, and successor problems''). \emph{Point of contact:} the I-state on a horizon pair is Israel's thermofield double \cite{Israel1976}, and its GNS representation is the Araki-Wyss doubling; the non-equilibrium state of the descent, the cosmological state problem, is outside this paper's scope (successor problem 5).
\end{remark}

\section{The mode continuum}

\phantomsection\label{c:note.hawking.setting-mode-continuum}
The boundary mode space of Definition \ref{c:obj.hawking.boundary-mode-algebra-windows-boost-dynamics} has countably many modes because the boundary modes, the ladder pairs of the blown-up singular boundary, one per point, are aggregated over the s-wave channels into the countable array of Theorem \ref{c:singularity.bosonisation-bridge-s-wave-sector}. The tower of iterated completions is the algebra's model, not the source of the count: its limit is the CAR algebra with the adjoined null pairs as modes (Fact \ref{c:hawking.imported-algebra-boundary-modes-car-algebra}), while on the physical closure one ascent completes the description (Corollary \emph{Answers to the iteration question} [T3, Cor 6.2]). The physical spectrum is emitted into a continuum of frequencies and a tower of angular channels, and this section carries the state theorem across that idealisation with the levels declared: what is exact, what is an inductive-limit result, and what is a modelling identification are separated.

\begin{definition}[Continuum mode space]\label{c:obj.hawking.continuum-mode-space}
Let $\mathfrak{h}$ be a separable complex Hilbert space (the one-particle space; for the Schwarzschild s-wave sector, $L^2(\R_+, d\omega)$; with angular channels, a direct sum of copies indexed by $(\ell, m)$ and species), and let $h$ be self-adjoint on $\mathfrak{h}$ (the one-particle Hamiltonian of Schwarzschild time; for the s-wave sector, multiplication by the frequency $\omega$; self-adjointness, not positivity, is what the theorems below use, and the section ``The finite boundary mode space, and the $\ell_P^2/A$ correction'''s bilateral ladder is an instance with spectrum unbounded below). The \emph{continuum mode space} is the CAR algebra $\CAR(\mathfrak{h})$ with the quasi-free (Bogoliubov) dynamics $\tau_t(a(f)) = a(e^{ith}f)$. The countable mode space of Definition \ref{c:obj.hawking.boundary-mode-algebra-windows-boost-dynamics} is the case $\mathfrak{h} = \ell^2(J)$; the CAR functor turns isometric inclusions of one-particle spaces into unital embeddings of algebras, so refining the tower's mode set into a dense subset of the continuum realises $\CAR(\mathfrak{h})$ as the completed union of tower-type subalgebras: the continuum is the tower's limit re-parametrised, not a new object.
\end{definition}

\begin{theorem}[The I-state on the continuum]\label{c:hawking.i-state-continuum}
\changedmark{gray}{reworded}
For the continuum mode space at $\beta = 2\pi/\kappa$, with $h$ self-adjoint:
\begin{enumerate}
\item the quasi-free state $\omega_I$ with symbol $A_I = (1 + e^{\beta h})^{-1}$ is $\tau$-KMS at $\beta$: the two-point KMS identity is the computation of clause (3) run in reverse on the core, and Fact \ref{c:hawking.imported-quasi-free-kms-reduces-two} lifts it from the two-point level to the state. When $h$ has pure point spectrum, the finite spectral windows (spans of finitely many eigenvectors) are literal instances of the section ``The state input: the I-state'''s exact window theorem, proved outright; this is the case of the tower's mode space, Definition \ref{c:obj.hawking.boundary-mode-algebra-windows-boost-dynamics};
\item the pairing of any quasi-free $\tau$-KMS state vanishes, by the computation of Theorem \ref{c:hawking.i-state}(2) performed on the core of clause (3);
\item the symbol derivation holds verbatim on the core: for $f$ in a bounded spectral subspace of $h$ (their union over cut-offs is a dense, $\tau$-invariant core on which $e^{-\beta h}f$ is defined and $a^\dagger(f)$ is entire analytic for $\tau$), the KMS identity forces $\ip{g}{A(1 + e^{-\beta h})f} = \ip{g}{e^{-\beta h}f}$, hence $A = A_I$ on the core and, $A$ and $A_I$ both bounded, everywhere: within the quasi-free class the continuum I-state is unique with derived two-point data. ($e^{-\beta h}$ may be unbounded when $h$ is unbounded below; the identity is only ever used on the core, and $A_I$ is bounded for every self-adjoint $h$.)
\end{enumerate}
Fact \ref{c:hawking.imported-ideal-fermi-gas-uniqueness} is stated for general self-adjoint $h$, so Corollary \ref{c:hawking.full-uniqueness-given-dynamics} also holds verbatim for the continuum mode space: among all states, $\omega_I$ is the unique $\tau$-KMS state.
\end{theorem}
\begin{proof}
(3) is the computation of Theorem \ref{c:hawking.i-state}(3) with the stated core in place of the window span: the core is $\tau$-invariant, $e^{-\beta h}$ preserves it and is bounded on it, and density plus boundedness of $A$ and $A_I$ extend the bounded identity $A = A_I$ from the core to $\mathfrak{h}$; the unbounded identity is used only on the core, where $1 + e^{-\beta h}$ acts bijectively. (2) is Theorem \ref{c:hawking.i-state}(2) on the same core. (1): the state $\omega_{A_I}$ exists for the symbol $0 \le A_I \le 1$ (Fact \ref{c:hawking.imported-quasi-free-states-car-algebra}); its two-point function satisfies the two-point KMS identity by (3) read backwards, and Fact \ref{c:hawking.imported-quasi-free-kms-reduces-two} upgrades the two-point identity to the KMS property of the state. What this paper certifies independently, and the script checks exactly, is the window case, where no imported calculus is needed. The closing uniqueness clause rests on \ref{c:hawking.imported-ideal-fermi-gas-uniqueness}, which gives a unique $\beta$-KMS state, the quasi-free state with symbol $(1 + e^{\beta h})^{-1}$, for any self-adjoint $h$.

\end{proof}

\begin{fact}[Imported: the thermal representation]\label{c:hawking.imported-thermal-representation}
The GNS representation of $\CAR(\mathfrak{h})$ in the state $\omega_I$ is the Araki-Wyss representation: a doubled Fock space carrying two commuting copies of the CAR, with the GNS vector entangling them pairwise and the modular operator of the GNS vector implementing exactly the flow $\tau$ at inverse temperature $\beta$. \emph{Provenance:} Araki-Wyss \cite{ArakiWyss1964}; textbook form \cite[\S 5.3]{BratteliRobinson1997}.
\emph{Source:} H. Araki and W. Wyss, Representations of canonical anticommutation relations, Helv. Phys. Acta 37 (1964), 136--159.; O. Bratteli and D. W. Robinson, Operator Algebras and Quantum Statistical Mechanics 2, 2nd ed., Springer, Berlin (1997)., \S 5.3.
\end{fact}

\begin{remark}[Free-fermion thermal field theory, and Israel's double]\label{c:note.hawking.free-fermion-thermal-field-theory-israels}
Fact \ref{c:hawking.imported-thermal-representation} is the continuum instance of T5's single-mode picture: the doubling is the completed pair structure mode by mode, the GNS vector is the thermofield double, and the modular flow is the boost. ``Free-fermion thermal field theory as the algebra carrying the mode continuum'' is now a theorem-backed description: the algebra is $\CAR(\mathfrak{h})$ (Fact \ref{c:hawking.imported-algebra-boundary-modes-car-algebra} plus Definition \ref{c:obj.hawking.continuum-mode-space}), the state is the I-state (Theorem \ref{c:hawking.i-state-continuum}), and the representation it generates is the standard thermal one. \emph{Point of contact:} Israel's thermofield reading of the eternal black hole \cite{Israel1976} is exactly this representation with the second copy read as the far wedge; T5's eviction dynamics runs inside it.
\end{remark}

\begin{remark}[Which dynamics, and at what level]\label{c:note.hawking.which-dynamics-what-level}
The one-particle Hamiltonian $h$ was supplied by the geometry: the Schwarzschild frequency of each boundary mode, with the temperature's coefficient fixed by the surface gravity (Fact \ref{c:hawking.imported-modular-descent-premises-geometric-normalisation}); in rapidity the weights are $\omega/\kappa$ at modular temperature $2\pi$. Two notes on scope. First, the greybody filtering of the section ``The four-dimensional lift'' does \emph{not} deform $h$ or the KMS structure: transmission through the exterior barrier is a property of the propagation from the horizon to infinity, not of the horizon algebra; conflating the two is what the interface premise (I) of Assembly \ref{c:hawking.greybody-inheritance-under-interface-premise} guards against, and it is named in that assembly's statement. Second, the identification of the mode set (which $(\omega, \ell, m)$ channels, with which degeneracies) is a modelling input from the standard near-horizon reduction; the theorems of this section are exact for whatever separable $(\mathfrak{h}, h)$ that input delivers; the algebra supplies the theorems and the physics supplies the input.
\end{remark}

\section{The four-dimensional lift}

\phantomsection\label{c:note.hawking.setting-four-dimensional-lift}
The results so far are stated for the boundary modes, in the effectively two-dimensional near-horizon theory whose unit is the hyperbolic atom. The physical spectrum is four-dimensional radiation sorted by helicity and angular momentum and filtered by the exterior geometry. This section assembles the lift. It is the paper's principal assembly section, its Tier B, and the helicity assembly's joints are itemised by name so that each carries its own provenance: one joint is a classical theorem, one is a theorem of this series completed in Proposition \ref{c:hawking.homogeneity-ordering-shift} below, and one is modelling, flagged as such; the greybody assembly itemises three more and adds one named premise, (I), its only non-theorem joint. What the lift does and does not claim is bounded in Remark \ref{c:note.hawking.what-lift-does-does-not-claim}.

\begin{fact}[Imported: the Penrose transform]\label{c:hawking.imported-penrose-transform}
\changedmark{gray}{reworded}
Massless free fields of helicity $s$ on (suitably convex regions of complexified) Minkowski space (complexified with the unit $I_4$, points $u + I_4 w$) are in natural isomorphism with the first cohomology of the corresponding region of the conformal spinor space $\CS$ (twistor space) with coefficients in $\mathcal{O}(-2s-2)$: the homogeneity $-2s-2$ classifies the helicity-$s$ field. \emph{Provenance:} Penrose \cite{Penrose1969}; the sheaf-cohomological form \cite{EPW1981}.
\emph{Source:} R. Penrose, Solutions of the zero-rest-mass equations, J. Math. Phys. 10 (1969), 38--39.; M. G. Eastwood, R. Penrose and R. O. Wells, Cohomology and massless fields, Commun. Math. Phys. 78 (1981), 305--351..
\end{fact}

\begin{proposition}[The homogeneity $-2s-2$ is the ordering shift]\label{c:hawking.homogeneity-ordering-shift}
\changedmark{gray}{reworded}
On functions on the conformal spinor space $\CS$ (twistor functions), realise the palatial relation of Fact \ref{c:hawking.imported-twistor-storey-ordering-shift} by $\bar Z_\nu = -\hbar\,\partial/\partial Z^\nu$ (the realisation with the palatial sign $[Z^\mu, \bar Z_\nu] = +\hbar\,\delta^\mu_\nu$), let $\Upsilon := Z^\mu \partial_\mu$ be the homogeneity (Euler) operator, its eigenvalue written $\mathsf{h}$ in sans serif to keep it clear of the one-particle Hamiltonian $h$, and let the helicity operator be one half the distinguished charge, $\mathsf{s} := \tfrac12\, q(\id)$, the standard helicity operator on $\CS$ (twistor helicity) \cite{Penrose1969, EPW1981} in the symmetric ordering that Fact \ref{c:hawking.imported-twistor-storey-ordering-shift} forces storey-wide. Then, as an operator identity,
\[
q(\id) \;=\; -\hbar\,(\Upsilon + 2),
\]
so a function on $\CS$ homogeneous of degree $\mathsf{h}$ is an eigenfunction of $\mathsf{s}$ with eigenvalue $s\hbar$ if and only if
\[
\mathsf{h} \;=\; -2s - 2.
\]
The $-2$ is exactly the ordering shift: without the symmetric-ordering shift of $-\hbar/2$ per complex dimension (over the complex unit $I_6$) the offset would vanish. Combined with Fact \ref{c:hawking.imported-penrose-transform}, the classification parameter of the Penrose transform is the completion's own $\mathfrak{u}(1)$ charge, its offset supplied by the ordering. The computation is not new to the series: it appears, without being stated as a result, inside the proof of \cite[Thm.~5.7]{ButlerT9}, where T9's script also checks it; what this proposition adds is the statement at result level, in the transform's context, discharging the homogeneity half of the successor problem T9 assigns to this paper (the transform itself stays open).
\end{proposition}
\begin{proof}
With $\bar Z_\nu = -\hbar\,\partial_\nu$,
\[
q(\id) \;=\; \tfrac12\sum_\mu\bigl(Z^\mu \bar Z_\mu + \bar Z_\mu Z^\mu\bigr)
\;=\; \sum_\mu Z^\mu \bar Z_\mu + \tfrac12\sum_\mu [\bar Z_\mu, Z^\mu]
\;=\; -\hbar\,\Upsilon - \tfrac12\cdot 4\hbar \;=\; -\hbar\,(\Upsilon + 2),
\]
the commutator term being $\sum_\mu[\bar Z_\mu, Z^\mu] = -4\hbar$ by the palatial relation, one $-\hbar$ per complex dimension of $\CS$ (over $I_6$); halved by the symmetrisation, this is Fact \ref{c:hawking.imported-twistor-storey-ordering-shift}'s shift of $-\hbar/2$ per complex dimension, totalling $-2\hbar$, and it is the only place the ordering enters. On a homogeneity-$\mathsf{h}$ eigenfunction, $\Upsilon f = \mathsf{h} f$, so $\mathsf{s} f = -\tfrac{\hbar}{2}(\mathsf{h} + 2)\,f$, and $\mathsf{s} f = s\hbar f$ if and only if $\mathsf{h} = -2s - 2$. The script verifies the displayed operator identity on all monomials of degree at most three in the four components of $Z \in \CS$, and the resulting homogeneity table at $s = \pm 1, \pm 2$. The reading of $\mathsf{h}$ as the classification parameter of the transform rests on \ref{c:hawking.imported-penrose-transform}, which classifies helicity-$s$ fields by homogeneity $-2s-2$.

\end{proof}

\begin{theorem}[The helicity lift, with its joints itemised]\label{c:hawking.helicity-lift-joints-itemised}
\changedmark{gray}{reworded}
Helicity is not additional structure imported into the framework: it is the eigenvalue of the distinguished $\mathfrak{u}(1)$ charge of the completion's conformal storey (the twistor storey: $\Cl(2,4)$ with its conformal spinor space $\CS$), and the four-dimensional mode space of each massless species has a canonical place in it. The joints, each with its provenance:
\begin{enumerate}
\item (\emph{Joint J1: the transform, classical.}) By Fact \ref{c:hawking.imported-penrose-transform}, the helicity-$s$ mode space is the cohomology at homogeneity $-2s-2$; photon kinematics are at homogeneities $-4$ and $0$ ($s = \pm1$), graviton kinematics at $-6$ and $+2$ ($s = \pm2$).
\item (\emph{Joint J2: the charge, proved in the series.}) The ordering shift is proved in the correspondence paper (Fact \ref{c:hawking.imported-twistor-storey-ordering-shift}), and Proposition \ref{c:hawking.homogeneity-ordering-shift} derives the homogeneity $-2s-2$ from it: the classification parameter of Joint J1 is an internal charge of the completion, so the helicity ladder is part of the ascended kinematics rather than a phenomenological label.
\item (\emph{Joint J3: the reduction, modelling.}) The current boson of the s-wave theory (Fact \ref{c:hawking.imported-s-wave-bosonic-sector}) is the s-wave reduction of each such four-dimensional field, and the lift from the boundary modes to the field is the standard partial-wave reconstruction, each helicity sector identified through the transform of Joint J1: the near-horizon reduction of field theory on the Schwarzschild exterior identifies each $(\omega, \ell, m, s)$ channel with a two-dimensional mode, which is where Definition \ref{c:obj.hawking.continuum-mode-space}'s modelling input enters, flagged as such in Remark \ref{c:note.hawking.which-dynamics-what-level}.
\end{enumerate}
The composite statement is the conjunction of the three joints with the provenance displayed: J1 is classical, J2 is Fact \ref{c:hawking.imported-twistor-storey-ordering-shift} together with Proposition \ref{c:hawking.homogeneity-ordering-shift}, and J3 is the one modelling joint. The homogeneity table is checked by the script.
\end{theorem}
\begin{proof}
The assembly itemises its joints in its statement; each joint carries its own provenance there.

Helicity is not additional structure imported into the framework: it is the eigenvalue of the distinguished $\mathfrak{u}(1)$ charge of the completion's conformal storey (the twistor storey: $\Cl(2,4)$ with its conformal spinor space $\CS$), and the four-dimensional mode space of each massless species has a canonical place in it. The joints, each with its provenance:
\begin{enumerate}
\item (\emph{Joint J1: the transform, classical.}) By Fact \ref{c:hawking.imported-penrose-transform}, the helicity-$s$ mode space is the cohomology at homogeneity $-2s-2$; photon kinematics are at homogeneities $-4$ and $0$ ($s = \pm1$), graviton kinematics at $-6$ and $+2$ ($s = \pm2$).
\item (\emph{Joint J2: the charge, proved in the series.}) The ordering shift is proved in the correspondence paper (Fact \ref{c:hawking.imported-twistor-storey-ordering-shift}), and Proposition \ref{c:hawking.homogeneity-ordering-shift} derives the homogeneity $-2s-2$ from it: the classification parameter of Joint J1 is an internal charge of the completion, so the helicity ladder is part of the ascended kinematics rather than a phenomenological label.
\item (\emph{Joint J3 (\ref{c:hawking.joint-j3-s-wave-reduction}): the reduction, modelling.}) The current boson of the s-wave theory (Fact \ref{c:hawking.imported-s-wave-bosonic-sector}) is the s-wave reduction of each such four-dimensional field, and the lift from the boundary modes to the field is the standard partial-wave reconstruction, each helicity sector identified through the transform of Joint J1: the near-horizon reduction of field theory on the Schwarzschild exterior identifies each $(\omega, \ell, m, s)$ channel with a two-dimensional mode, which is where Definition \ref{c:obj.hawking.continuum-mode-space}'s modelling input enters, flagged as such in Remark \ref{c:note.hawking.which-dynamics-what-level}.
\end{enumerate}
The composite statement is the conjunction of the three joints with the provenance displayed: J1 is classical, J2 is Fact \ref{c:hawking.imported-twistor-storey-ordering-shift} together with Proposition \ref{c:hawking.homogeneity-ordering-shift}, and J3 is the one modelling joint. The homogeneity table is checked by the script.

\end{proof}

\begin{theorem}[Greybody inheritance, under the interface premise]\label{c:hawking.greybody-inheritance-under-interface-premise}
\changedmark{gray}{reworded}
\emph{Premise (I), the interface identification:} the horizon-sector state feeds the exterior modes in the standard way, i.e.\ the occupancy of each near-horizon channel is carried to the corresponding asymptotic channel by the exterior propagation. This is the same assumption every derivation of Hawking flux makes, from \cite{Hawking1975} onward, and it is this assembly's one non-theorem joint. Under (I), the emitted flux in each channel $(\omega, \ell, m, s)$ equals (with $\mathcal{N}$ the count of emitted quanta)
\[
\frac{d\mathcal{N}}{dt\,d\omega} \;=\; \frac{1}{2\pi}\,\Gamma_{s\ell m}(\omega)\; n_\pm(\omega), \qquad
n_\pm(\omega) = \frac{1}{e^{\omega/\TH} \pm 1}
\]
(the upper sign for the fermionic species, the lower for the bosonic sector), with $\Gamma_{s\ell m}$ the classical transmission coefficient through the Regge-Wheeler or Teukolsky barrier. The joints: (i) the barriers and their transmission coefficients are properties of the exterior metric, and by Fact \ref{c:hawking.imported-exterior-exactly-general-relativity} the framework's exterior description \emph{is} classical general relativity on the invertible region, exactly and not asymptotically, so $\Gamma_{s\ell m}$ is inherited verbatim, mode by mode \cite{ReggeWheeler1957, Teukolsky1973}; (ii) the occupancies at the horizon are Theorem \ref{c:hawking.i-state-continuum} for the fermionic species and Fact \ref{c:hawking.imported-s-wave-bosonic-sector}'s (Q)-conditional half for the bosonic sector, both at exactly $\TH$ with the temperature's coefficient fixed geometrically (Fact \ref{c:hawking.imported-modular-descent-premises-geometric-normalisation}); (iii) given (I), the displayed formula is the definition of filtered thermal flux: the framework's flux prediction equals Hawking's \cite{Hawking1975, Page1976}, greybody structure included.
\end{theorem}
\begin{proof}
The assembly itemises its joints in its statement; each joint carries its own provenance there.

\emph{Premise (I) (\ref{c:hawking.premise-i-interface-identification}), the interface identification:} the horizon-sector state feeds the exterior modes in the standard way, i.e.\ the occupancy of each near-horizon channel is carried to the corresponding asymptotic channel by the exterior propagation. This is the same assumption every derivation of Hawking flux makes, from \cite{Hawking1975} onward, and it is this assembly's one non-theorem joint. Under (I), the emitted flux in each channel $(\omega, \ell, m, s)$ equals (with $\mathcal{N}$ the count of emitted quanta)
\[
\frac{d\mathcal{N}}{dt\,d\omega} \;=\; \frac{1}{2\pi}\,\Gamma_{s\ell m}(\omega)\; n_\pm(\omega), \qquad
n_\pm(\omega) = \frac{1}{e^{\omega/\TH} \pm 1}
\]
(the upper sign for the fermionic species, the lower for the bosonic sector), with $\Gamma_{s\ell m}$ the classical transmission coefficient through the Regge-Wheeler or Teukolsky barrier. The joints: (i) the barriers and their transmission coefficients are properties of the exterior metric, and by Fact \ref{c:hawking.imported-exterior-exactly-general-relativity} the framework's exterior description \emph{is} classical general relativity on the invertible region, exactly and not asymptotically, so $\Gamma_{s\ell m}$ is inherited verbatim, mode by mode \cite{ReggeWheeler1957, Teukolsky1973}; (ii) the occupancies at the horizon are Theorem \ref{c:hawking.i-state-continuum} for the fermionic species and Fact \ref{c:hawking.imported-s-wave-bosonic-sector}'s (Q)-conditional half for the bosonic sector, both at exactly $\TH$ with the temperature's coefficient fixed geometrically (Fact \ref{c:hawking.imported-modular-descent-premises-geometric-normalisation}); (iii) given (I), the displayed formula is the definition of filtered thermal flux: the framework's flux prediction equals Hawking's \cite{Hawking1975, Page1976}, greybody structure included.

\end{proof}

\begin{remark}[What the lift does and does not claim]\label{c:note.hawking.what-lift-does-does-not-claim}
The section's claim is deliberately bounded: the \emph{spectrum} of the four-dimensional radiation is derived (occupancies theorem-grade at the horizon under the named premises, filtering inherited exactly on the invertible region, the helicity location derived from the ordering shift). What is not claimed is a derivation of four-dimensional field \emph{dynamics} from the boundary modes: that is the task of the action paper T10 \cite{ButlerT10} and, for the emergent transverse geometry that would make Joint J3 itself a theorem, of the entropy paper T8 \cite{ButlerT8}, the fuzzy-sphere question being shared between T8 and T10. The lift beyond assembly level is a named successor problem (the section ``Scope, and successor problems'').
\end{remark}

\section{The finite boundary mode space, and the $\ell_P^2/A$ correction}

\phantomsection\label{c:note.hawking.setting-finite-boundary-mode-space-p}
The bosonisation bridge of T5 rests, at one joint, on the charge-zero Fock isomorphism of the \emph{infinite} chiral fermion \cite{FrenkelKac1980}, while the physical boundary mode space is finite: $L$ modes, with $L$ of order the horizon area $A$ in Planck units. This section proves that everything the finite window changes is supported at its edges, with the exact operator identity and quantified tails, and then records the reading: the corrections are area-suppressed, so the joint's finite-mode correction takes the form of a prediction.

\subsection*{The exact edge identity}
\phantomsection\label{c:note.hawking.setting-exact-edge-identity}
Fix a window of $L$ fermionic modes $\psi_1, \dots, \psi_L$ (CAR: $\{\psi_k, \psi_l^\dagger\} = \delta_{kl}$, all other anticommutators zero) and the currents
\[
a_m \;:=\; \sum_{k=1}^{L-m} \psi_k^\dagger\,\psi_{k+m}, \qquad 1 \le m < L .
\]

\begin{proposition}[The current commutator, exactly, with its edge split]\label{c:hawking.current-commutator-exactly-edge-split}
In the finite window, exactly and as an operator identity,
\[
[\,a_m,\, a_m^\dagger\,] \;=\; B_m - T_m, \qquad
B_m := \sum_{k=1}^{m} N_k, \qquad T_m := \sum_{k=L-m+1}^{L} N_k ,
\]
with $N_k = \psi_k^\dagger\psi_k$: the deviation from the infinite-volume central value is a difference of \emph{edge operators}, supported on the lowest and highest $m$ modes. Consequently, for $2m \le L$, on every state in which the bottom $m$ levels are filled and the top $m$ levels are empty (an \emph{edge-clean} state), $[a_m, a_m^\dagger]$ acts as the scalar $m$: the free-boson relation holds exactly, not asymptotically, on the edge-clean subspace. Moreover $[a_m, a_{m'}] = 0$ for all $m, m'$, exactly; and for $m \neq m'$ the cross commutator $[a_m, a_{m'}^\dagger]$ is a difference of bilinears supported on the lowest and highest $\max(m, m')$ modes, which annihilates every state that is edge-clean at depth $\max(m, m')$.
\end{proposition}
\begin{proof}
The fermionic bilinear commutator is $[\psi_a^\dagger\psi_b,\ \psi_c^\dagger\psi_d] = \delta_{bc}\,\psi_a^\dagger\psi_d - \delta_{ad}\,\psi_c^\dagger\psi_b$, an instance of the CAR (\ref{c:std-algebra.witt-basis-car}). With $a_m^\dagger = \sum_{l=1}^{L-m}\psi_{l+m}^\dagger\psi_l$:
\[
[a_m, a_m^\dagger] = \sum_{k,l=1}^{L-m}\bigl(\delta_{k+m,\,l+m}\,\psi_k^\dagger\psi_l - \delta_{k l}\,\psi_{l+m}^\dagger\psi_{k+m}\bigr)
= \sum_{k=1}^{L-m} N_k - \sum_{k=m+1}^{L} N_k ,
\]
after shifting the index of the second sum. The shared middle range cancels pairwise: for $2m \le L$ the terms $k = m{+}1, \dots, L{-}m$ appear in both sums and cancel, leaving $B_m - T_m$; for $2m > L$ the two sums instead share the empty middle and the overlap $k = L{-}m{+}1, \dots, m$ of $B_m$ and $T_m$ cancels on the right-hand side, so the identity $\sum_{1}^{L-m} N_k - \sum_{m+1}^{L} N_k = B_m - T_m$ holds in both regimes. On an edge-clean state (which requires $2m \le L$ to be non-vacuous), $B_m$ acts as $m$ (each bottom mode occupied) and $T_m$ as $0$. For the cross commutators, the same computation gives
\[
[a_m, a_{m'}] = \sum_{k}\psi_k^\dagger\psi_{k+m+m'} - \sum_{k}\psi_k^\dagger\psi_{k+m+m'} = 0
\]
exactly, both surviving sums running over the same range $1 \le k \le L - m - m'$; and, for $m > m'$,
\[
[a_m, a_{m'}^\dagger] = \sum_{k=1}^{m'}\psi_k^\dagger\psi_{k+m-m'} \;-\; \sum_{k=L-m+1}^{L-m+m'}\psi_k^\dagger\psi_{k+m-m'},
\]
every surviving index within $\max(m,m')$ of an edge. Each bottom term annihilates an edge-clean state because its target level $k$ is occupied; each top term because its source level $k + m - m'$ lies in the empty top band. The case $m < m'$ follows by taking adjoints, $[a_m, a_{m'}^\dagger] = [a_{m'}, a_m^\dagger]^\dagger$, with the same band support; it annihilates edge-clean states by the mirrored argument (creation onto an occupied bottom level, lowering of an empty top level). The script verifies the identity on the full Fock space of an eight-mode window, state by state.

\end{proof}

\subsection*{The Planck factor, and the window tails}
\phantomsection\label{c:note.hawking.setting-planck-factor-window-tails}
This subsection works on the bilateral ladder, $k \in \Z$ with equal level spacing $\epsilon$: the idealisation of a deep sea, an instance of Definition \ref{c:obj.hawking.continuum-mode-space} with $h = \epsilon\,\mathrm{diag}(k)_{k \in \Z}$, unbounded below. The currents of the section ``The exact edge identity'' extend to the bilateral sums $a_m = \sum_{k \in \Z}\psi_k^\dagger\psi_{k+m}$, read term by term inside I-state expectations, where the sums converge absolutely; the finite window $[-K, K]$ returns in clause (3).

\begin{proposition}[Bose-Einstein occupancy, with quantified window corrections]\label{c:hawking.bose-einstein-occupancy-quantified-window-corrections}
Let the ladder carry the linear s-wave weight assignment (the convention T5's Premise (Q) adopts, a modelling input here as there; Fact \ref{c:hawking.imported-s-wave-bosonic-sector}): level spacing $\epsilon$, sea at level $0$, I-state occupancies $n_k = (x^{-k} + 1)^{-1}$ with $x = e^{-\beta\epsilon} \in (0,1)$ the Boltzmann factor of the level spacing. Then:
\begin{enumerate}
\item \emph{(Detailed-balance identity.)} For every $k$ and $m \ge 1$, exactly,
\[
n_k - n_{k+m} \;=\; \bigl(x^{-m} - 1\bigr)\; n_{k+m}\,\bigl(1 - n_k\bigr).
\]
\item \emph{(The Planck factor, exactly, in the bilateral limit.)} Summing (1) over $k \in \Z$,
\[
\sum_{k \in \Z} n_{k+m}\,(1 - n_k) \;=\; \frac{m}{x^{-m} - 1} \;=\; \frac{m}{e^{\beta\epsilon m} - 1}:
\]
the current two-point function $\langle a_m^\dagger a_m\rangle_I = \sum_k n_{k+m}(1 - n_k)$ (a Wick computation) is exactly $m$ times the Bose-Einstein value at $\TH$, with no vacuum correction in this normalisation. The factor $m$ is the free-boson normalisation: on the bilateral ladder the currents are used only inside expectations, and $\langle [a_m, a_m^\dagger]\rangle_I = \sum_k (n_k - n_{k+m}) = m$, the telescoped left side of (1), matching the finite-window scalar of Proposition \ref{c:hawking.current-commutator-exactly-edge-split}; the unit-normalised boson is $\hat a_m = a_m/\sqrt{m}$, whose occupancy $\langle \hat a_m^\dagger \hat a_m\rangle_I = 1/(e^{\beta\epsilon m} - 1)$ is the $n_-$ of the lift.
\item \emph{(Window tails.)} For the window $[-K, K]$ with the sea at the centre, i.e.\ the sum of (2) restricted to $k \in [-K, K-m]$ (both mode indices inside the window), the bilateral sum exceeds the window sum by a positive amount bounded by $2m\,x^{K - m}/(1 - x)$: exponentially small in the distance from the sea to the edges, in units of the level spacing.
\end{enumerate}
\end{proposition}
\begin{proof}
(1) With $u := x^{-k}$ and $w := x^{-m}$ (the I-state occupancies of \ref{c:hawking.i-state} on the linear s-wave weights of \ref{c:hawking.imported-s-wave-bosonic-sector}): $n_k - n_{k+m} = \tfrac{1}{u+1} - \tfrac{1}{uw+1} = \tfrac{u(w-1)}{(u+1)(uw+1)}$, and $(w-1)\,n_{k+m}(1-n_k) = (w-1)\cdot\tfrac{1}{uw+1}\cdot\tfrac{u}{u+1}$: identical. The script checks it as an identity of rational functions in the formal variable $u$ at several rational $w$.
(2) The terms $n_{k+m}(1 - n_k)$ decay like $x^{|k|}$ in both directions, so all sums converge absolutely and may be regrouped. Fix a residue class $k = r + jm$, $j \in \Z$: the class sum of the left side of (1) is $\lim_{j_0 \to \infty}\sum_{j=-j_0}^{j_0}\bigl(n_{r+jm} - n_{r+(j+1)m}\bigr) = \lim_{j_0 \to \infty}\bigl(n_{r-j_0 m} - n_{r+(j_0+1)m}\bigr) = 1 - 0$, the interior terms cancelling in adjacent pairs; there are $m$ residue classes, and dividing by $(x^{-m}-1)$ via (1) gives the display. The Wick computation: in the gauge-invariant quasi-free state, $\langle \psi^\dagger_{l+m}\psi_l\,\psi^\dagger_k\psi_{k+m}\rangle$ has vanishing disconnected part for $m \ge 1$ (the two-point function is mode-diagonal) and connected part $\delta_{kl}\,n_{k+m}(1 - n_k)$.
(3) Bound each term two ways, both valid for every integer index: $n_{k+m}(1-n_k) \le n_{k+m} \le x^{k+m}$ and $n_{k+m}(1-n_k) \le 1 - n_k \le x^{-k}$, since $n_j \le x^j$ and $1 - n_j \le x^{-j}$ are each the statement $1 \le 1 + x^{\pm j}$. The terms missing from the window sum are those with $k > K - m$ or $k < -K$, all positive, so the difference is positive and at most
\[
\sum_{k > K-m} x^{k+m} + \sum_{k < -K} x^{-k} \;=\; \frac{2\,x^{K+1}}{1-x} \;\le\; \frac{2m\,x^{K-m}}{1-x}.
\]
The script verifies instances of the bound exactly, and checks that the window sums at $K = 30$ and $K = 40$ agree with the closed form to within the proven tails, all in exact rational arithmetic.

\end{proof}

\begin{remark}[The $\ell_P^2/A$ correction, interpreted]\label{c:note.hawking.p-2-correction-interpreted}
Propositions \ref{c:hawking.current-commutator-exactly-edge-split} and \ref{c:hawking.bose-einstein-occupancy-quantified-window-corrections} together are the quantitative form of T5's note on the level of its assembly joint. The charge-zero Fock isomorphism is exact for the infinite chiral fermion; for the physical boundary mode space of $L \sim A/\ell_P^2$ modes, every deviation is edge-supported (Proposition \ref{c:hawking.current-commutator-exactly-edge-split}), and thermal expectations of the deviations carry explicit tails (Proposition \ref{c:hawking.bose-einstein-occupancy-quantified-window-corrections}(3)). The corrections therefore vanish as the area grows, at a rate set by the window geometry: the relative scale of generic finite-size effects is read dimensionally as $1/L = \ell_P^2/A$ (for the discretisation of mode sums and the Sugawara normalisation; proving that scaling per named observable is successor problem 4), while edge-clean bulk correlators at fixed $m$ are exponentially smaller, at the proven rates. The quantum-gravity corrections are area-suppressed and vanish where general relativity is well tested; any new physics in this derivation would appear in the $\ell_P^2/A$ correction, and that correction is a prediction of the derivation rather than a defect in it.
\end{remark}

\begin{remark}[Point of contact: the Weinberg-Witten triage]\label{c:note.hawking.point-contact-weinberg-witten-triage}
The bosonic sector's lift contains a graviton candidate (homogeneities $-6, +2$ in Assembly \ref{c:hawking.helicity-lift-joints-itemised}), and a composite graviton must answer \cite{WeinbergWitten1980}: a theory with a Lorentz-covariant conserved stress tensor admits no massless spin-2 particle coupling to it. The evasion, recorded here as this paper's own claim and at Tier C: the theorem's hypotheses require the stress tensor and the spin-2 state to be defined on one and the same spacetime, and on the composite reading they are not; the covariant stress tensor of the boundary modes is defined in the two-dimensional reduced theory at the degenerate locus, the graviton candidate is a collective mode of the emergent four-dimensional description, and the hypotheses fail exactly where they fail for boundary gravitons in holography. The evidential triage, stated because the series states its own claims at their level: by \cite{Weinberg1964}, \emph{any} massless spin-2, composite or fundamental, couples universally to stress-energy, so the equivalence principle and the universal vertex are common ground and do not discriminate; what is distinctive to the composite reading is the ultraviolet behaviour and the $\ell_P^2/A$ correction of Remark \ref{c:note.hawking.p-2-correction-interpreted}, and the discriminating experiments, if any, are at the scale of that correction. Below the Planck scale no available experiment discriminates the composite reading from the fundamental one, so the reading is at present untestable; this paper says so rather than presenting the absence of a refutation as support. The engineering no-go and its ontology-independence are assigned to the action paper T10 \cite{ButlerT10}.
\end{remark}

\section{Scope, and successor problems}

\subsection*{What is proved, what is imported, what is modelling}
\phantomsection\label{c:note.hawking.setting-what-proved-what-imported-what}
The abstract states each claim at one of three tiers: Tier A for what this paper proves, Tier B for what it assembles from results proved elsewhere or from named premises, and Tier C for an interpretation offered as such. The first paragraph below is the Tier A content, Proposition \ref{c:hawking.homogeneity-ordering-shift} included, which the abstract meets as the lift's one proved joint; the second lists the imports that content rests on, certified by their source papers or the literature; the third is the assembly and modelling content, the assembly at Tier B and the modelling items stated as such, with the composite-graviton triage at Tier C. In the claim table of T17 \cite{ButlerT17} the paper as a whole is Tier B, because its physical conclusions pass through the assembly section, while its mathematical core is Tier A.

\emph{Proved in this paper, at the standard of \cite{ButlerT1}:} the single-mode modular structure (Proposition \ref{c:hawking.modular-structure-completed-pair}); the Gibbs/KMS equivalence with uniqueness in finite dimensions and the restriction lemma (Lemmas \ref{c:hawking.gibbs-states-are-kms-conversely-finite}, \ref{c:hawking.kms-restricts-invariant-subalgebras}). Also proved here: the I-state theorem's window-exact existence (Theorem \ref{c:hawking.i-state}(1), window case), the forced gauge invariance, and the derivation of the Fermi-Dirac symbol from the KMS identity, on the boundary modes and on the continuum core (Theorems \ref{c:hawking.i-state}(2)--(3), \ref{c:hawking.i-state-continuum}(2)--(3)); the homogeneity $-2s-2$ from the ordering shift (Proposition \ref{c:hawking.homogeneity-ordering-shift}); the exact edge identity, the edge-clean exactness of the free-boson relation, and the cross-commutator support bound (Proposition \ref{c:hawking.current-commutator-exactly-edge-split}); the detailed-balance identity, the exact bilateral Planck factor, and the window tail bounds (Proposition \ref{c:hawking.bose-einstein-occupancy-quantified-window-corrections}, on the (Q)-convention ladder named in its statement).

\emph{Imported as Facts with provenance:} T5's modular descent under its premises (P) and (R), with the geometric normalisation (Fact \ref{c:hawking.imported-modular-descent-premises-geometric-normalisation}); the tower's CAR limit (Fact \ref{c:hawking.imported-algebra-boundary-modes-car-algebra}); the s-wave bosonisation results, their occupancy half under T5's premise (Q) (Fact \ref{c:hawking.imported-s-wave-bosonic-sector}); the exterior exactness at $\Lambda = 0$ (Fact \ref{c:hawking.imported-exterior-exactly-general-relativity}); the twistor storey and the ordering shift (Fact \ref{c:hawking.imported-twistor-storey-ordering-shift}). From the literature: Araki's quasi-free calculus, the inductive-limit KMS result, the two-point reduction, and the ideal-Fermi-gas uniqueness (Facts \ref{c:hawking.imported-quasi-free-states-car-algebra}, \ref{c:hawking.imported-kms-passes-inductive-limits}, \ref{c:hawking.imported-quasi-free-kms-reduces-two}, \ref{c:hawking.imported-ideal-fermi-gas-uniqueness}); the Araki-Wyss representation (Fact \ref{c:hawking.imported-thermal-representation}); the Penrose transform (Fact \ref{c:hawking.imported-penrose-transform}); the soft-graviton universality \cite{Weinberg1964} used in Remark \ref{c:note.hawking.point-contact-weinberg-witten-triage}.

\emph{Assembly and modelling, flagged at their level:} the helicity lift's joint J3 and the greybody interface premise (I) (the section ``The four-dimensional lift''); the identification of the physical mode set, and the linear s-wave weight convention of the bilateral ladder (Remark \ref{c:note.hawking.which-dynamics-what-level}, Proposition \ref{c:hawking.bose-einstein-occupancy-quantified-window-corrections}); the $\ell_P^2/A$ dimensional reading of generic finite-size corrections (Remark \ref{c:note.hawking.p-2-correction-interpreted}); the composite-graviton reading and its triage (Remark \ref{c:note.hawking.point-contact-weinberg-witten-triage}), which is Tier C and says so.

\subsection*{Successor problems}
\phantomsection\label{c:note.hawking.setting-successor-problems}
Six problems for the series' list of open directions.

\begin{enumerate}
\item \emph{State selection from geometry.} Formulate the geometric-modular-action axiom for a net of boundary-mode algebras over the horizon and derive the boost dynamics, the temperature entering through the net's geometric normalisation rather than as an algebra-intrinsic quantity, which T5's rate-rigidity theorem shows it cannot be \cite[Thm.~7.16]{ButlerT5}. Proposition \ref{c:hawking.modular-structure-completed-pair}(5) fixes the single-atom orbits, and T5's factor-level analysis (its condition (G), rejected as over-determining) records what is already established.
\item \emph{The lift beyond assembly.} Derive joint J3, the partial-wave identification, from the emergent transverse geometry (the fuzzy sphere of the counting programme), converting the helicity lift (Assembly \ref{c:hawking.helicity-lift-joints-itemised}) from assembly to theorem; shared with the entropy paper T8 \cite{ButlerT8} and the action paper T10 \cite{ButlerT10}.
\item \emph{The Kerr continuum.} Superradiant channels make the rotating analogue of Theorem \ref{c:hawking.i-state-continuum} genuinely different (no global KMS state at fixed angular velocity for all modes); formulate the correct rotating statement.
\item \emph{The $\ell_P^2/A$ correction, per observable.} Refine Remark \ref{c:note.hawking.p-2-correction-interpreted} into a table of leading finite-$L$ corrections for named observables (spectrum distortion, correlator deficits), with the $\ell_P^2/A$ scaling proved rather than dimensional.
\item \emph{The cosmological state.} Connect the equilibrium I-state to the non-equilibrium state of the descent.
\item \emph{Helicity counting per boundary mode.} Decide whether one boundary mode of the near-horizon theory carries one helicity channel of the lift or both helicities of a massless species. The two answers differ downstream by a factor of two in the capacity per boundary mode ($\ln 2$ against $2\ln 2$) and so in the entropy paper's calibration residue, $\kappa_c = 1/(4\ln 2)$ against $1/(8\ln 2)$ \cite[Prop.~4.12]{ButlerT8}. In this paper the question sits at joint J3 of Assembly \ref{c:hawking.helicity-lift-joints-itemised}, the partial-wave identification of each $(\omega, \ell, m, s)$ channel with a two-dimensional mode, which is modelling; joints J1 and J2 place the two helicities at distinct homogeneities but do not say how many channels one boundary mode registers. T8 takes the one-helicity reading, at the tier of that joint; settling it is part of converting J3 to a theorem (problem 2).
\end{enumerate}

\section{Position in the series}

\phantomsection\label{c:note.hawking.setting-position-series}
This paper is T7 of the series. \emph{It relies on six predecessors:} the primer T0 \cite{ButlerT0} and the root paper T1 \cite{ButlerT1} for the completed-pair language and the transport principle; the tower paper T3 \cite{ButlerT3} for the CAR limit that makes the boundary modes an algebra (Fact \ref{c:hawking.imported-algebra-boundary-modes-car-algebra}); the singularity paper T5 \cite{ButlerT5} for modular descent under its premises (P) and (R), the bosonisation bridge under (Q), and the surface gravity $\kappa = 1/4M$ (Facts \ref{c:hawking.imported-modular-descent-premises-geometric-normalisation} and \ref{c:hawking.imported-s-wave-bosonic-sector}); the interior paper T6 \cite{ButlerT6} for the exterior exactness that greybody inheritance rests on (Fact \ref{c:hawking.imported-exterior-exactly-general-relativity}); and the twistor paper T9 \cite{ButlerT9} for the ordering shift from which Proposition \ref{c:hawking.homogeneity-ordering-shift} is promoted.

\emph{Papers that rely on this one:} the entropy paper T8 \cite{ButlerT8} imports the single-mode modular structure, the quasi-free calculus, and the finite-window identities; the action paper T10 \cite{ButlerT10} imports the finite-window quantification of the current algebra; the standard-model paper T12 \cite{ButlerT12} imports the helicity identification of Proposition \ref{c:hawking.homogeneity-ordering-shift}; and the capstone T17 \cite{ButlerT17} tabulates this paper's claims at their tiers.

\section{Machine verification}

\phantomsection\label{c:note.hawking.setting-machine-verification}
The computational identities asserted in this paper are verified by the script deposited with it, \texttt{verify\_T07.py} \cite{ButlerVerifyT7}. The script imports the series' shared exact-arithmetic library \texttt{GA.py}, runs entirely in exact arithmetic (rational numbers and rational functions in a formal Boltzmann variable, no floating point anywhere), and so certifies each check as an identity, or as an exactly evaluated inequality, rather than an approximation. Twenty-five checks pass. Not scripted, and stated as such where they occur: the imported Facts, certified by their source papers' scripts or by the literature; the continuum passage beyond its window content; the assembly joints of the section ``The four-dimensional lift''; the classical Penrose transform; and the geometric inputs $\kappa = 1/4M$ and $\TH = \kappa/2\pi$, proved and scripted in T5.

The coverage, section by section.

\begin{description}
\item[the section ``Setting and imports'', the single mode (Proposition \ref{c:hawking.modular-structure-completed-pair}).] All five clauses in the basis-blade arithmetic over $\Cl(1,1)$: the null pair's squares and anticommutator, $\sigma = e_0e_1$ with $\sigma^2 = 1$, the one-sided weight relations, the finite conjugation at the rational rotor $R = \tfrac53 + \tfrac43\,\sigma$ (that is, $\cosh(\lambda/2) = \tfrac53$) with scale factors $\tfrac19$ and $9$, the projector identities, and the one-dimensionality of the bivector space; plus the Fock identity $\tfrac12\sigma = N - \tfrac12$.
\item[the section ``The state input: the I-state'', the state (Lemma \ref{c:hawking.gibbs-states-are-kms-conversely-finite}, Theorem \ref{c:hawking.i-state}).] The Gibbs KMS identity on all sixteen matrix-unit pairs of the mode algebra, as identities of rational functions in the formal Boltzmann factor; the uniqueness solve, performed symbolically, the KMS constraints on an arbitrary density matrix having a one-dimensional kernel spanned by the Gibbs matrix. Then the symbol itself: the Fermi-Dirac detailed-balance identity $n/(1-n) = e^{-\beta\omega}$; the Fermi-Dirac product state on a three-mode window confirmed equal to the Gibbs matrix of $H_F$, the two sides built independently, then confirmed KMS on all $4096$ matrix-unit pairs; and, on a two-mode window, the full KMS constraint system solved for an arbitrary density matrix against independently supplied dynamics, with one-dimensional kernel, vanishing pairing, and the Fermi-Dirac symbol.
\item[the section ``The mode continuum'', the continuum.] No separately computable identity: its exact content is the window case, scripted under Section 3, and its remaining steps are the imported calculus of Facts \ref{c:hawking.imported-quasi-free-kms-reduces-two} and \ref{c:hawking.imported-ideal-fermi-gas-uniqueness}, certified by their sources.
\item[the section ``The four-dimensional lift'', the lift (Proposition \ref{c:hawking.homogeneity-ordering-shift}).] The operator identity $q(\id) = -\hbar(\Upsilon + 2)$ on all twistor monomials of degree at most three, and the homogeneity table $-2s-2 \mapsto (-4, 0, -6, +2)$ at $s = \pm 1, \pm 2$.
\item[the section ``The finite boundary mode space, and the $\ell_P^2/A$ correction'', the finite window (Propositions \ref{c:hawking.current-commutator-exactly-edge-split} and \ref{c:hawking.bose-einstein-occupancy-quantified-window-corrections}).] The edge identity $[a_m, a_m^\dagger] = B_m - T_m$ as an operator identity on all $256$ states of an eight-mode window for $m = 1, 2, 3$; the exact free-boson relation on every edge-clean state; the exact vanishing of $[a_m, a_{m'}]$ and the vanishing of the cross commutator on edge-clean states. Then the ladder: the detailed-balance identity of Proposition \ref{c:hawking.bose-einstein-occupancy-quantified-window-corrections}(1) in the formal $x^{-k}$ at three values of $x^{-m}$, with the Bose-Einstein form $\bar n/(1 + \bar n) = e^{-\beta\omega}$ of the clause-(2) occupancy, at $\omega = \epsilon m$ on the ladder; the window sums' approach to the exact Planck factor from below, within the proven tail bound, as exact rational inequalities at two temperatures and three $m$; and the Wick formula checked against the direct density-matrix expectation on the full $256$-state window, the ladder in the paper's occupancy convention.
\end{description}

The script exits nonzero on any failure and prints one line per check, tagged by the statement it certifies.

\section{Results used from elsewhere}
The proofs above use the following results of the series, stated briefly; each is proved in the paper named.

\begin{fact}[The Clifford side of the ordering shift]\label{c:correspondence.clifford-side-ordering-shift}
In the Fock model the grading element of an adjoined null pair acts on spinors as $N - \tfrac12$, with $N$ the occupation number of the corresponding mode.
\emph{Source:} T9, Prop 5.6 (\href{https://doi.org/10.5281/zenodo.22722290}{doi:10.5281/zenodo.22722290}).
\end{fact}

\begin{fact}[The Fock model of the complexified ascended algebra]\label{c:correspondence.fock-model-complexified-ascended-algebra}
The exterior algebra on a maximal isotropic subspace, with wedge and contraction operators, is the unique irreducible module of the complexified Clifford algebra $\Cl(2,4)_\C$, and its even and odd halves are the two four-dimensional chiralities.
\emph{Source:} T9, Thm 3.1 (\href{https://doi.org/10.5281/zenodo.22722290}{doi:10.5281/zenodo.22722290}).
\end{fact}

\begin{fact}[The shared core, and the ordering dictionary]\label{c:correspondence.shared-core-ordering-dictionary}
Under the Klein identification, the palatial algebra's symmetrised quadratics and the Clifford bivectors share the same $\mathfrak{so}(2,4)$ core acting compatibly on conformal spinors, and both storeys carry the same signed ordering shift of minus one half per dimension.
\emph{Source:} T9, Thm 5.7 (\href{https://doi.org/10.5281/zenodo.22722290}{doi:10.5281/zenodo.22722290}).
\end{fact}

\begin{fact}[Symmetrised quadratics of the palatial algebra]\label{c:correspondence.symmetrised-quadratics-palatial-algebra}
The symmetrised quadratics of the palatial algebra realise $\mathfrak{gl}(4,\C)$ faithfully through commutators and equal the normal-ordered ones minus the central term $\tfrac{\hbar}{2}\operatorname{tr}M$, a shift of $-2\hbar$ for the helicity generator.
\emph{Source:} T9, Lem 5.5 (\href{https://doi.org/10.5281/zenodo.22722290}{doi:10.5281/zenodo.22722290}).
\end{fact}

\begin{fact}[Covariant form of the vacuum equations]\label{c:interior.covariant-form-vacuum-equations}
On the invertible region of a gravitational field on flat spacetime, the vacuum field equations hold exactly when the null scale tractor $\nInf$ of the derived metric is parallel for the tractor gauge field $\Omega_{T}$.
\emph{Source:} T6, Thm 4.5 (\href{https://doi.org/10.5281/zenodo.22721990}{doi:10.5281/zenodo.22721990}).
\end{fact}

\begin{fact}[Imported: equivalence with the Einstein equations]\label{c:interior.imported-equivalence-einstein-equations}
Wherever $\hb$ is invertible, the GTG field equations are equivalent to Einstein's equations for the induced metric, and in vacuum to its Ricci-flatness.
\emph{Source:} A. Lasenby, C. Doran and S. Gull, Gravity, gauge theories and geometric algebra, Phil. Trans. R. Soc. Lond. A 356 (1998), 487--582., \S 4 (\S 4.3, covariant forms of the field equations), App.~C; C. Doran and A. Lasenby, Geometric Algebra for Physicists, Cambridge University Press (2003)., Ch.~13-14.
\end{fact}

\begin{fact}[The bosonisation bridge, s-wave sector]\label{c:singularity.bosonisation-bridge-s-wave-sector}
Number-conserving bilinears of the boundary fermion modes satisfy free-boson commutation relations exactly, up to a precisely identified edge term on finite windows, and under Premise (Q) they carry Bose-Einstein occupancy at $\TH$ in the s-wave sector.
\emph{Source:} T5, Thm 8.2 (\href{https://doi.org/10.5281/zenodo.22721790}{doi:10.5281/zenodo.22721790}).
\end{fact}

\begin{fact}[$\kappa = 1/4M$, computed]\label{c:singularity.computed}
For the Schwarzschild field on flat spacetime the redshifted acceleration of static observers, $v = K/\sqrt f$, is $f'(r)/2$, so the surface gravity is $\kappa = 1/4M$.
\emph{Source:} T5, Prop A.2 (\href{https://doi.org/10.5281/zenodo.22721790}{doi:10.5281/zenodo.22721790}).
\end{fact}

\begin{fact}[The descended form is state independent]\label{c:singularity.descended-form-state-independent}
Whatever the ascended state, each descended single-mode marginal is a Gibbs state of the boost generator at a unique $\beta_k$, for every mode of nonzero weight.
\emph{Source:} T5, Thm 7.6 (\href{https://doi.org/10.5281/zenodo.22721790}{doi:10.5281/zenodo.22721790}).
\end{fact}

\begin{fact}[The Euclidean section, and its period]\label{c:singularity.euclidean-section-period}
The Euclidean section of the Schwarzschild field is smooth at its tip, the Euclidean counterpart of the bifurcation sphere, if and only if imaginary Killing time has period $p = 2\pi/\kappa = 8\pi M$; any other period leaves a conical point of total angle $\kappa p$.
\emph{Source:} T5, Prop 7.18 (\href{https://doi.org/10.5281/zenodo.22721790}{doi:10.5281/zenodo.22721790}).
\end{fact}

\begin{fact}[The matching: Steps 4b and 4c join by equation]\label{c:singularity.matching-steps-4b-4c-join-equation}
The modular bivector of the completed pair, $\sigma = e_0e_1$, satisfies $S(\sigma) = -2\sigma$ and $\sigma^2 = +1$, which characterise the curvature-selected boost plane up to sign, so the modular flow and the curvature select the same plane, identified at $r > 0$ and transported to the singular boundary.
\emph{Source:} T5, Cor 6.4 (\href{https://doi.org/10.5281/zenodo.22721790}{doi:10.5281/zenodo.22721790}).
\end{fact}

\begin{fact}[Modular descent]\label{c:singularity.modular-descent}
A pure, boost-invariant, entangling state of the horizon pair descends to a KMS state of the phase flow $\operatorname{Ad}(e^{it\sigma/2})$ at a free $\beta$, and under the geometric normalisation (R) to the Fermi-Dirac state at exactly $\TH = \kappa/2\pi = 1/8\pi M$.
\emph{Source:} T5, Thm 7.3 (\href{https://doi.org/10.5281/zenodo.22721790}{doi:10.5281/zenodo.22721790}).
\end{fact}

\begin{fact}[The rate, and the discharge of hypothesis (iii)]\label{c:singularity.rate-discharge-hypothesis-iii}
Under (P) and (R) every mode's marginal is the Gibbs state at the Hawking temperature $T = 1/8\pi M$, so each marginal is faithful and the cyclic-and-separating hypothesis of modular descent is derived rather than assumed.
\emph{Source:} T5, Thm 7.20 (\href{https://doi.org/10.5281/zenodo.22721790}{doi:10.5281/zenodo.22721790}).
\end{fact}

\begin{fact}[Rate rigidity: the temperature is not algebra intrinsic]\label{c:singularity.rate-rigidity-temperature-not-algebra-intrinsic}
Automorphisms of the single-mode algebra that intertwine the boost's phase flow $\operatorname{Ad}(e^{itK})$ move its KMS states only by $\beta \mapsto \pm\beta$, so the value of the temperature depends on a parametrisation of the flow that the algebra does not supply.
\emph{Source:} T5, Thm 7.16 (\href{https://doi.org/10.5281/zenodo.22721790}{doi:10.5281/zenodo.22721790}).
\end{fact}

\begin{fact}[Exponentials of bivectors with scalar square]\label{c:std-algebra.bivector-exponential}
If $A^{2} = \lambda \in \R$ then $e^{A}$ is $\cos\theta + \frac{\sin\theta}{\theta}A$, $\cosh\theta + \frac{\sinh\theta}{\theta}A$ or $1 + A$, a rotor with $e^{A}\widetilde{e^{A}} = 1$.
\emph{Source:} the algebra core of the Transport Lab, proved there with an exact check; standard, see C. Doran and A. Lasenby, Geometric Algebra for Physicists, Cambridge University Press (2003), Ch. 4 (foundations of geometric algebra: rotors, versors).
\end{fact}

\begin{fact}[The product of two vectors]\label{c:std-algebra.vector-squares}
$v^{2} = Q(v)$ and $uv + vu = 2B(u, v)$, so $uv = B(u, v) + \tfrac12(uv - vu)$.
\emph{Source:} the algebra core of the Transport Lab, proved there with an exact check; standard, see C. Doran and A. Lasenby, Geometric Algebra for Physicists, Cambridge University Press (2003), Ch. 4 (foundations of geometric algebra).
\end{fact}

\begin{fact}[Null pairs satisfy the canonical anticommutation relations]\label{c:std-algebra.witt-basis-car}
For a Witt basis ($B(v_i,v_j) = B(w_i,w_j) = 0$, $B(v_i,w_j) = \tfrac12\delta_{ij}$), $\{v_i,v_j\} = \{w_i,w_j\} = 0$ and $\{v_i,w_j\} = \delta_{ij}$ in $\Cl(V,Q)$.
\emph{Source:} the algebra core of the Transport Lab, proved there with an exact check; standard, see C. Doran and A. Lasenby, Geometric Algebra for Physicists, Cambridge University Press (2003), Ch. 9 (multiparticle states).
\end{fact}

\begin{fact}[The algebraic limit, even charge]\label{c:tower.algebraic-limit-even-charge}
For even charge the algebraic direct limit of the tower is canonically $\Cl(p,q) \otimes \Mat_{2^\infty}$, equivalently the Clifford algebra of $\R^{p,q} \perp H_1 \perp H_2 \perp \cdots$.
\emph{Source:} T3, Thm 5.3 (\href{https://doi.org/10.5281/zenodo.22721556}{doi:10.5281/zenodo.22721556}).
\end{fact}

\begin{fact}[The CAR limit, even charge]\label{c:tower.car-limit-even-charge}
For $p - q$ even, the C*-completion of the complexified tower limit is $\Cl(p,q)_\C \otimes \CAR$, with the adjoined null vectors as the fermionic mode operators.
\emph{Source:} T3, Thm 5.6 (\href{https://doi.org/10.5281/zenodo.22721556}{doi:10.5281/zenodo.22721556}).
\end{fact}

\begin{fact}[The ladder pairs satisfy the CAR, as vectors]\label{c:tower.ladder-pairs-satisfy-car-vectors}
In the $k$-th tower algebra the $k$ adjoined null pairs satisfy the canonical anticommutation relations of $k$ fermionic modes exactly, as vectors, and each new rung adds one mode without altering the existing relations.
\emph{Source:} T3, Prop 4.1 (\href{https://doi.org/10.5281/zenodo.22721556}{doi:10.5281/zenodo.22721556}).
\end{fact}

\begin{prediction}[The no-echo prediction, registered]\label{c:scrambling.no-echo-prediction-registered}
Under Postulate (C1), which puts the quantum structure at the stretched singularity, post-merger ringdown is the unmodified Kerr quasinormal spectrum, with no echo train at the delay $8M\ln(M/m_P)$.
\emph{Kill condition:} One statistically confirmed echo train at the Planck-structure delay.
\end{prediction}

\section*{Changes since the deposited version}
\addcontentsline{toc}{section}{Changes since the deposited version}
Since the deposited version 1.0: 1 result corrected, 2 changed, 9 reworded (wording, notation or citations only) and 26 cited results added. The comparison is with the source of version 1.0 on Zenodo, \href{https://doi.org/10.5281/zenodo.22722121}{doi:10.5281/zenodo.22722121}. A result counts as revised when the words or mathematics of its statement or proof differ from that source; a citation added in a parenthesis does not count. A result the paper cites from elsewhere is compared with the deposit of the paper that proves it. Each revision is classed as corrected, changed or reworded; the classes were assigned by a review of each difference and the store's record of why it was made, and a revision not yet reviewed is counted as changed.
\subsection*{Changed}
The content changed by a later decision, not because the deposited version was wrong.
\begin{itemize}\setlength{\itemsep}{2pt}
\item Assembled~\ref{c:hawking.imported-modular-descent-premises-geometric-normalisation} (\emph{Imported: modular descent, its premises, and the geometric normalisation}), Fact 2.2 in the deposit: statement and proof. The flow for which the descended state is KMS is now named as the phase flow of the same generator; provenance pointers added. \emph{Store notes.} \emph{2026-10-21:} The two boost flows: flow\_note recorded; formal, proof main, technical now name the flow (original formal in source\_formal, proof in orig-proof-*.tex where changed; version bumped).
\item Definition~\ref{c:obj.hawking.boundary-mode-algebra-windows-boost-dynamics} (\emph{The boundary-mode algebra, windows, and the boost dynamics}), Def 2.7 in the deposit: statement. New provenance for the mode set and a new convention placing modes below the level by particle-hole conjugates, zero mode excluded. \emph{Store notes.} \emph{2026-10-21:} A definition cannot carry uses or assumes; the added links are about edges from the claims whose statements use the defined notion. \emph{2026-10-21:} The two boost flows: flow\_note recorded.
\end{itemize}
\subsection*{Reworded}
Wording, notation, terminology or citations only; the mathematics and the claim are as deposited.
\begin{itemize}\setlength{\itemsep}{2pt}
\item Proposition~\ref{c:hawking.modular-structure-completed-pair} (\emph{Modular structure of the completed pair}), Prop 2.1 in the deposit: proof. Citations of standard algebra facts and cross-references only.
\item Assembled~\ref{c:hawking.imported-algebra-boundary-modes-car-algebra} (\emph{Imported: the algebra of the boundary modes is the CAR algebra}), Fact 2.3 in the deposit: statement. Citations and cross-references only; statement repeated in the proof slot.
\item Assembled~\ref{c:hawking.imported-s-wave-bosonic-sector} (\emph{Imported: the s-wave bosonic sector}), Fact 2.4 in the deposit: proof. Citations and cross-references only; statement repeated in the proof slot.
\item Assembled~\ref{c:hawking.imported-twistor-storey-ordering-shift} (\emph{Imported: the conformal storey (twistor storey), and the ordering shift}), Fact 2.6 in the deposit: statement and proof. Wording only: twistor space written as the conformal spinor space over I\_6; the ordering shift is unchanged.
\item Theorem~\ref{c:hawking.i-state-continuum} (\emph{The I-state on the continuum}), Thm 4.2 in the deposit: proof. Cross-references added, with a sentence naming the source of the uniqueness clause.
\item Fact~\ref{c:hawking.imported-penrose-transform} (\emph{Imported: the Penrose transform}), Fact 5.1 in the deposit: statement. Wording only: twistor-space region written as a region of the conformal spinor space, complex unit named I\_4.
\item Proposition~\ref{c:hawking.homogeneity-ordering-shift} (\emph{The homogeneity $-2s-2$ is the ordering shift}), Prop 5.2 in the deposit: statement and proof. Wording only: twistor functions written as functions on the conformal spinor space; computation unchanged.
\item Theorem~\ref{c:hawking.helicity-lift-joints-itemised} (\emph{The helicity lift, with its joints itemised}), Asm 5.3 in the deposit: statement and proof. Wording only: twistor storey called the conformal storey; joints and provenance unchanged.
\item Theorem~\ref{c:hawking.greybody-inheritance-under-interface-premise} (\emph{Greybody inheritance, under the interface premise}), Asm 5.4 in the deposit: proof. Citations and cross-references only; statement repeated in the proof slot.
\end{itemize}
\noindent In addition, 26 results from elsewhere in the series or the algebra core are stated here that the deposited paper cited or used without stating; 1 result gained citations only. These are not marked.
\begin{thebibliography}{99}
\bibitem{ButlerT5} L.~Butler,
\emph{Division by Zero at the Gravitational Singularity: The Transport Principle Applied to the Black Hole, and the Information Paradox as a Category Error},
Transport Papers T5 (2026), doi:\href{https://doi.org/10.5281/zenodo.22721790}{10.5281/zenodo.22721790}.
\bibitem{ButlerT3} L.~Butler,
\emph{The Tower and the CAR Limit: Iterated Transport, the Frozen Clock, and Second Quantisation},
Transport Papers T3 (2026), doi:\href{https://doi.org/10.5281/zenodo.22721556}{10.5281/zenodo.22721556}.
\bibitem{ButlerT9} L.~Butler,
\emph{The Twistor Correspondence: Ascended Spinors, Incidence, and the Palatial Bridge},
Transport Papers T9 (2026).
\bibitem{ButlerT6} L.~Butler,
\emph{The Geometry of the Black-Hole Interior: Where Division Fails, the Weighted Limit, and the Kerr Tests},
Transport Papers T6 (2026), doi:\href{https://doi.org/10.5281/zenodo.22721990}{10.5281/zenodo.22721990}.
\bibitem{Hawking1975} S.~W.~Hawking,
Particle creation by black holes,
\emph{Commun.\ Math.\ Phys.} \textbf{43} (1975), 199--220.
\bibitem{Israel1976} W.~Israel,
Thermo-field dynamics of black holes,
\emph{Phys.\ Lett.\ A} \textbf{57} (1976), 107--110.
\bibitem{BisognanoWichmann1976} J.~J.~Bisognano and E.~H.~Wichmann,
On the duality condition for quantum fields,
\emph{J.\ Math.\ Phys.} \textbf{17} (1976), 303--321.
\bibitem{Sewell1982} G.~L.~Sewell,
Quantum fields on manifolds: PCT and gravitationally induced thermal states,
\emph{Ann.\ Phys.} \textbf{141} (1982), 201--224.
\bibitem{Araki1971} H.~Araki,
On quasifree states of CAR and Bogoliubov automorphisms,
\emph{Publ.\ RIMS Kyoto Univ.} \textbf{6} (1971), 385--442.
\bibitem{BratteliRobinson1997} O.~Bratteli and D.~W.~Robinson,
\emph{Operator Algebras and Quantum Statistical Mechanics 2}, 2nd ed.,
Springer, Berlin (1997).
\bibitem{BDFS2000} D.~Buchholz, O.~Dreyer, M.~Florig and S.~J.~Summers,
Geometric modular action and spacetime symmetry groups,
\emph{Rev.\ Math.\ Phys.} \textbf{12} (2000), 475--560.
\bibitem{Penrose1969} R.~Penrose,
Solutions of the zero-rest-mass equations,
\emph{J.\ Math.\ Phys.} \textbf{10} (1969), 38--39.
\bibitem{EPW1981} M.~G.~Eastwood, R.~Penrose and R.~O.~Wells,
Cohomology and massless fields,
\emph{Commun.\ Math.\ Phys.} \textbf{78} (1981), 305--351.
\bibitem{ReggeWheeler1957} T.~Regge and J.~A.~Wheeler,
Stability of a Schwarzschild singularity,
\emph{Phys.\ Rev.} \textbf{108} (1957), 1063--1069.
\bibitem{Teukolsky1973} S.~A.~Teukolsky,
Perturbations of a rotating black hole I,
\emph{Astrophys.\ J.} \textbf{185} (1973), 635--647.
\bibitem{Page1976} D.~N.~Page,
Particle emission rates from a black hole,
\emph{Phys.\ Rev.\ D} \textbf{13} (1976), 198--206.
\bibitem{FrenkelKac1980} I.~B.~Frenkel and V.~G.~Kac,
Basic representations of affine Lie algebras and dual resonance models,
\emph{Invent.\ Math.} \textbf{62} (1980), 23--66.
\bibitem{WeinbergWitten1980} S.~Weinberg and E.~Witten,
Limits on massless particles,
\emph{Phys.\ Lett.\ B} \textbf{96} (1980), 59--62.
\bibitem{Weinberg1964} S.~Weinberg,
Photons and gravitons in S-matrix theory,
\emph{Phys.\ Rev.} \textbf{135} (1964), B1049--B1056.
\bibitem{ButlerT1} L.~Butler,
\emph{The Transport Principle: Total Division in Clifford Algebras via Minimal Completion},
Transport Papers T1 (2026).
\bibitem{ButlerVerifyT7} L.~Butler,
\emph{verify\_T07.py: exact-arithmetic verification of the computational identities of ``Hawking Radiation from the Ledger''},
Python script, deposited with this paper alongside the shared library GA.py (2026).
\bibitem{ArakiWyss1964} H.~Araki and W.~Wyss,
Representations of canonical anticommutation relations,
\emph{Helv.\ Phys.\ Acta} \textbf{37} (1964), 136--159.
\bibitem{ButlerT10} L.~Butler,
\emph{Dark Energy and the Action of Gravity: MacDowell-Mansouri, the Frozen Point, and One Way to Grip},
Transport Papers T10 (2026), doi:\href{https://doi.org/10.5281/zenodo.22722496}{10.5281/zenodo.22722496}.
\bibitem{ButlerT8} L.~Butler,
\emph{Black-Hole Entropy and Scrambling: The Singularity Ledger and the No-Echo Prediction},
Transport Papers T8 (2026), doi:\href{https://doi.org/10.5281/zenodo.22722176}{10.5281/zenodo.22722176}.
\bibitem{ButlerT17} L.~Butler,
\emph{The Transport Programme: Method, Status, and the Falsifiable Ledger},
Transport Papers T17 (2026).
\bibitem{ButlerT0} L.~Butler,
\emph{An Introduction to Geometric Algebra: A Primer for the Transport Papers},
Transport Papers T0 (2026).
\bibitem{ButlerT12} L.~Butler,
\emph{The Standard Model from the Minimal Massive Composite},
Transport Papers T12 (2026), doi:\href{https://doi.org/10.5281/zenodo.22722669}{10.5281/zenodo.22722669}.
\end{thebibliography}
\end{document}